\documentclass[11pt]{article}
\usepackage[a4paper,margin=21mm]{geometry}
\usepackage{fontspec}
\setmainfont{Latin Modern Roman}
\newfontfamily\CJKfont{Noto Serif CJK SC}
\XeTeXlinebreaklocale "zh"
\XeTeXlinebreakskip=0pt plus 1pt
\usepackage{amsmath,amssymb,amsthm,amscd,graphicx,color}
\usepackage{xurl}
\usepackage[unicode,hidelinks]{hyperref}
\allowdisplaybreaks
\setlength\emergencystretch{3em}
\numberwithin{equation}{section}

\usepackage{bm}
\usepackage{eucal}

\numberwithin{equation}{section}

\newtheorem{Theorem}{Theorem}[section]
\newtheorem{Lemma}[Theorem]{Lemma}
\newtheorem{Proposition}[Theorem]{Proposition}
 { \theoremstyle{definition}
\newtheorem{Definition}[Theorem]{Definition}
\newtheorem{Example}[Theorem]{Example}
\newtheorem{Remark}[Theorem]{Remark} }

\newcommand{\pr}[1]{#1^{\prime}}
\newcommand{\pf}{\operatorname{Pf}}
\newcommand{\tr}{{\operatorname{Tr}}}
\newcommand{\mfrak}[1]{\mathfrak{#1}}
\newcommand{\mcal}[1]{\mathcal{#1}}
\newcommand{\mbb}[1]{\mathbb{#1}}
\newcommand{\mrm}[1]{\mathrm{#1}}
\newcommand{\what}[1]{\widehat{#1}}

\begin{document}
\begin{center}\Large Finite occupied-and-empty conditioning preserves projection Pfaffian point processes under the stated doubled-space regularity condition\end{center}
\noindent\textbf{Stable ID:} 1339\quad\textbf{Author:} Shinji Koshida\par
\noindent\textbf{Work:} Pfaffian Point Processes from Free Fermion Algebras: Perfectness and Conditional Measures\par
\noindent\textbf{Source:} \url{https://sigma-journal.com/2021/008/}\par
\noindent\textbf{Version:} Official SIGMA 17 (2021) article 008, DOI 10.3842/SIGMA.2021.008; journal PDF and native TeX acquired 2026-09-30, exact bytes pinned by SHA256\par
\noindent\textbf{Rights:} CC BY 4.0. Authors retain copyright. \url{https://creativecommons.org/licenses/by/4.0/}. Reuse and typesetting adaptation; original language and mathematical content preserved.\par
\noindent\textbf{Difficulty (prior selection preserved):} {\CJKfont H2: The proof reduces finite mixed conditioning to single-point steps while tracking the projection regularity and the doubled-space involution. It derives two successive projection reductions, computes the conditioned Pfaffian correlations, and identifies the resulting kernel with the projection onto the explicitly modified subspace. This simultaneously verifies the probability-law and Hilbert-space descriptions, rather than just applying a finite Pfaffian Schur complement.}\par
\medskip\noindent\textbf{Editorial boundary note (not author text):} Selected complete Theorem1.17, finite disjoint occupied/empty sets and exact doubled-space (X,X-prime)-regularity, with full original projection/subspace formula. Full PfPP correlation/kernel definitions, CAR/quasi-free correspondence and author proof, Fock basis and existence/Lenard positivity proof are retained, followed by every conditional projection, finite decomposition, regularity, two-step matrix reduction, occupied and empty correlation computation and final identification. Author symmetric finite-set/single-point cases explicitly described as similar remain literal bounded repetitions with both formulas written, without an editor expansion. Standard prior CAR representation and finite projection background stay precise references. The projection reduction is explicitly positioned as new; perfectness/Schur/continuous-system/future quasi-invariance outcomes are unselected. Every printed index/inner-product variant remains unchanged.\par
\medskip\hrule\medskip\noindent\textbf{Author text begins below.}\par
\setcounter{section}{1}\setcounter{subsection}{0}\setcounter{Theorem}{0}\setcounter{equation}{0}
% Original source lines 60--196
\subsection{Pfaffian point process}
In this paper, we assume that $\mfrak{X}$ is a countable set. The collection of point configurations in~$\mfrak{X}$ is identified with $\Omega=\Omega (\mfrak{X})=\{0,1\}^{\mfrak{X}}$, which is equipped with the product topology to be a~compact topological space. We regard each element $\omega\in \Omega$ as a function $\omega\colon \mfrak{X}\to\{0,1\}$ or a~collection of points $\omega=\{x_{i}\in\mfrak{X}\}_{i}$. We adopt the $\sigma$-algebra of Borel sets $\Sigma$. Then, for distinct points $x_{1},\dots, x_{n}\in \mfrak{X}$, the cylinder set
\begin{equation*}
	\Omega_{x_{1},\dots, x_{n}}=\{\omega\in\Omega\,|\,\omega (x_{1})=\cdots =\omega (x_{n})=1\}
\end{equation*}
is measurable.
Given a probability measure $M$ on $(\Omega,\Sigma)$, the $n$-point correlation function $\rho^{M}_{n}$, $n\in\mbb{N}$ is an $n$-variable symmetric function defined as the probability weight of cylinder sets:
\begin{equation*}
	\rho^{M}_{n}(x_{1},\dots, x_{n})=M(\Omega_{x_{1},\dots, x_{n}}),
\end{equation*}
where $x_{1},\dots, x_{n}\in \mfrak{X}$ are distinct. It is conventional to extend $\rho^{M}_{n}$ to a function on $\mfrak{X}^{n}$ so that it vanishes if any two points coincide.
Note that a system of correlation functions $\big\{\rho^{M}_{n}\big\}_{n\in\mbb{N}}$ determines the probability measure $M$ uniquely since the cylinder sets generate $\Sigma$.
A random variable $X$ with values in $(\Omega,\Sigma)$ is called a point process in $\mfrak{X}$. We also call a probability measure on $(\Omega,\Sigma)$ a point process not distinguishing a random variable from its distribution.

To define a Pfaffian point process, we need to fix some notations.
Suppose that a $(2\times 2)$-matrix-valued function $\mbb{K}(\cdot,\cdot)\colon \mfrak{X}\times\mfrak{X}\to M(2;\mbb{C})$ satisfying the anti-symmetry
\begin{equation}\label{eq:anti-sym_matrix_func}
	\mbb{K}(x,y)^{\mrm{T}}=-\mbb{K}(y,x),\qquad x, y\in\mfrak{X}
\end{equation}
is given. For any $n\in\mbb{N}$, we adopt the following identification of vector spaces:
\begin{align*}
	M(2;\mbb{C})\otimes M(n;\mbb{C}) & \xrightarrow{\sim} M(2n;\mbb{C}),\\
	 e_{i,j}\otimes e_{k,l}&\mapsto e_{2(k-1)+i, 2(l-1)+j},\qquad i,j=1,2,\quad k,l=1,\dots, n,
\end{align*}
where $e_{i,j}$ is the matrix with entry $1$ at the $(i,j)$-position and $0$ elsewhere.
Given points $x_{1},\dots, x_{n}\in\mfrak{X}$, we understand $[\mbb{K}(x_{i},x_{j})]_{1\le i,j\le n}$ as a $(2n\times 2n)$-matrix under the above identification and by regarding $x_{1},\dots, x_{n}$ as legs for $M(n;\mbb{C})$. The anti-symmetry (\ref{eq:anti-sym_matrix_func}) of the function $\mbb{K}(\cdot,\cdot)$ ensures the anti-symmetry of the $(2n\times 2n)$-matrix $[\mbb{K}(x_{i},x_{j})]_{1\le i,j\le n}$ so that
\begin{equation*}
[\mbb{K}(x_{i},x_{j})]_{1\le i,j\le n}^{\mrm{T}}=\big[\mbb{K}(x_{j},x_{i})^{\mrm{T}}\big]_{1\le i,j\le n}=-[\mbb{K}(x_{i},x_{j})]_{1\le i,j\le n}.
\end{equation*}
Therefore, the Pfaffian $\pf [\mbb{K}(x_{i},x_{j})]_{1\le i,j\le n}$ is defined.

\begin{Definition}
A probability measure $M$ on $(\Omega,\Sigma)$ is a Pfaffian point process (PfPP) if there exists a $(2\times 2)$-matrix-valued function $\mbb{K}^{M}(\cdot,\cdot)\colon \mfrak{X}\times\mfrak{X}\to M(2;\mbb{C})$ satisfying the anti-symmetry~(\ref{eq:anti-sym_matrix_func}) such that every $n$-point function with $n\in\mbb{N}$ admits a Pfaffian expression
\begin{equation*}
	\rho_{n}^{M}(x_{1},\dots, x_{n})=\pf \big[\mbb{K}^{M}(x_{i},x_{j})\big]_{1\le i,j\le n}.
\end{equation*}
We call the matrix-valued function $\mbb{K}^{M}(\cdot,\cdot)$ a correlation kernel of the PfPP $M$.
\end{Definition}

For a PfPP $M$, we write its correlation kernel as
\begin{equation*}
	\mbb{K}^{M}(x,y)
	=\left(
	\begin{matrix}
		\mbb{K}_{11}^{M}(x,y) & \mbb{K}_{12}^{M}(x,y) \\
		\mbb{K}_{21}^{M}(x,y) & \mbb{K}_{22}^{M}(x,y)
	\end{matrix}
	\right),\qquad x,y\in\mfrak{X}.
\end{equation*}
If, in particular, $\mbb{K}_{11}^{M}(x,y)=\mbb{K}_{22}^{M}(x,y)=0$ identically holds, we have
\begin{equation*}
	\mrm{Pf}\big[\mbb{K}^{M}(x_{i},x_{j})\big]_{1\le i,j\le n}=\det \big[\mbb{K}_{12}^{M}(x_{i},x_{j})\big]_{1\le i,j\le n}.
\end{equation*}
Therefore, we may naturally have the following definition as a special case of PfPPs.

\begin{Definition}
A PfPP $M$ on $(\Omega,\Sigma)$ is called a determinantal point process (DPP) if $\mbb{K}^{M}_{11}(x,y)\allowbreak =\mbb{K}^{M}_{22}(x,y)=0$ identically.
In this case, each correlation function admits a determinantal expression
\begin{equation*}
	\rho_{n}^{M}(x_{1},\dots, x_{n})=\det \big[K^{M}(x_{i},x_{j})\big]_{1\le i,j\le n},\qquad K^{M}(x,y)=\mbb{K}^{M}_{12}(x,y),\qquad x,y\in\mfrak{X}
\end{equation*}
We adopt the abuse of terminology to call the function $K^{M}(\cdot,\cdot)$ on $\mfrak{X}\times\mfrak{X}$ a correlation kernel of the DPP $M$.
\end{Definition}

DPPs form a significant class of point processes that have many applications to random partition and asymptotic representation theory \cite{Borodin1998a, BorodinOlshanski1998a, Borodin1998b,BorodinOlshanski1998b,Okounkov2001,Olshanski1998, Olshanski2003} (see also \cite{Borodin2011,BenHoughKrishnapurPeresVirag2006, Lyons2003, Soshnikov2000}). The analogy between DPPs and fermionic calculi is well-known \cite{Lytvynov2002,LytvynovMei2007,Olshanski2020, ShiraiTakahashi2003a, ShiraiTakahashi2003b, Spohn1987}. In particular, a quasi-free state of a certain class over an algebra of anti-commutation relations gives a DPP.
Compared to DPPs, there seem to be less attempts to unify an interplay between PfPPs and free fermions, while there are many preceding studies on concrete examples (incomplete references include \cite{BorodinRains2005, Ferrari2004, KatoriTanemura2007,Matsumoto2005, MatsumotoShirai2013,Nagao2007,Rains2000,Vuletic2007,WangLi2019}).
One of the aims of this paper is to extend a portion of the insight on DPPs found in the above mentioned literatures and to establish a framework to study PfPPs from free fermionic perspectives. Many settings, notations and problems are motivated by \cite{Olshanski2020}.

\subsection{From a positive contraction to a PfPP}
A well-known construction of a DPP on $\mfrak{X}$ \cite{Soshnikov2000} starts from an operator $K$ on $\ell^{2}(\mfrak{X})$ such that $K=K^{\ast}$ and $0\le K\le 1$, namely, a~positive contraction on~$\ell^{2}(\mfrak{X})$. Given a~positive contraction~$K$ on $\ell^{2}(\mfrak{X})$, there exists a DPP $M^{K}$ on $\mfrak{X}$ such that each correlation function is given by
\begin{equation*}
	\rho^{M^{K}}_{n}(x_{1},\dots, x_{n})=\det [K(x_{i},x_{j}) ]_{1\le i,j\le n},\quad K(x,y)=(e_{x},Ke_{y})_{\ell^{2}(\mfrak{X})}.
\end{equation*}
Here $e_{x}\in \ell^{2}(\mfrak{X})$, $x\in\mfrak{X}$ is a function defined by $e_{x}(y)=\delta_{x,y}$, $y\in\mfrak{X}$. Then the collection $\{e_{x}\}_{x\in\mfrak{X}}$ forms a complete orthonormal system of $\ell^{2}(\mfrak{X})$.

Here, we give a direct generalization of this result to PfPPs.
We define the complex conjugate~$J$ of $\ell^{2}(\mfrak{X})$ as an anti-linear operator on it that fixes each function~$e_{x}$, $x\in\mfrak{X}$.
We set $\mcal{K}=\ell^{2}(\mfrak{X})\oplus \ell^{2}(\mfrak{X})$ and take an anti-unitary involution $\Gamma$ on $\mcal{K}$ defined by
\begin{equation}
\label{eq:Gamma}
	\Gamma=\left(
	\begin{matrix}
	0 & J \\
	J & 0
	\end{matrix}
	\right)
\end{equation}
according to the prescribed direct sum decomposition.
For this pair $(\mcal{K},\Gamma)$, we consider the collection of operators
\begin{equation}
\label{eq:collection_covariance_ops}
	\mcal{Q}(\mcal{K},\Gamma)=\big\{S\in\mfrak{B}(\mcal{K})\,|\,0\le S=S^{\ast}\le 1,\, \overline{S}=1-S\big\},
\end{equation}
where $\mfrak{B}(\mcal{K})$ is the set of bounded operators on $\mcal{K}$, and for any operator $A\in\mfrak{B}(\mcal{K})$, we write $\overline{A}:=\Gamma A\Gamma$.
We will show that each operator $S\in\mcal{Q}(\mcal{K},\Gamma)$ uniquely determines a PfPP.

\begin{Proposition}\label{prop:existence_ppp}
Let us take $S\in\mcal{Q}(\mcal{K},\Gamma)$ and write it as
\begin{equation*}
	S=\left(
	\begin{matrix}
		S_{11} & S_{12} \\
		S_{21} & S_{22}
	\end{matrix}
	\right).
\end{equation*}
There exists a PfPP $M^{S}$ on $\mfrak{X}$ such that each correlation function is given by
\begin{equation*}
	\rho^{M^{S}}_{n}(x_{1},\dots, x_{n})=\mrm{Pf} [\mbb{K}_{S}(x_{i},x_{j}) ]_{1\le i,j\le n},
\end{equation*}
where
\begin{equation}
\label{eq:matrix_kernel_S}
	\mbb{K}_{S}(x,y)=\left(
	\begin{matrix}
		(e_{x},S_{21}e_{y})_{\ell^{2}(\mfrak{X})} & (e_{x},S_{22}e_{y})_{\ell^{2}(\mfrak{X})} \\
		(e_{x},(S_{11}-1)e_{y})_{\ell^{2}(\mfrak{X})} & (e_{x},S_{12}e_{y})_{\ell^{2}(\mfrak{X})}
	\end{matrix}
	\right),\qquad x,y\in\mfrak{X}.
\end{equation}
\end{Proposition}

\begin{Remark}
If $S$ is diagonal; $S_{12}=S_{21}=0$, the PfPP $M^{S}$ is a DPP.
\end{Remark}

\begin{Remark}
\label{rem:check_anti-sym_matrix_func}
For $S\in\mcal{Q}(\mcal{K},\Gamma)$, the function $\mbb{K}_{S}(\cdot,\cdot)$ satisfies the required anti-symmetry (\ref{eq:anti-sym_matrix_func}). In fact, the self-adjointness of $S$ gives $S_{11}^{\ast}=S_{11}$, $S_{22}^{\ast}=S_{22}$, $S_{12}^{\ast}=S_{21}$, and the identity $\overline{S}=1-S$ implies $JS_{11}J=1-S_{22}$, $JS_{12}J=-S_{21}$. In particular, we have $S_{12}^{\mrm{T}}=-S_{12}$, $S_{21}^{\mrm{T}}=-S_{21}$. Therefore,
\begin{gather*}
	(e_{x},S_{21}e_{y})_{\ell^{2}(\mfrak{X})} =-(e_{y},S_{21}e_{x})_{\ell^{2}(\mfrak{X})},\\
	(e_{x},S_{12}e_{y})_{\ell^{2}(\mfrak{X})} =-(e_{y},S_{12}e_{x})_{\ell^{2}(\mfrak{X})}, \\
	(e_{x},S_{22}e_{y})_{\ell^{2}(\mfrak{X})} =\overline{(e_{y},J(1-S_{11})Je_{x})_{\ell^{2}(\mfrak{X})}}=-(e_{y},(S_{11}-1)e_{x})_{\ell^{2}(\mfrak{X})}, \\
(e_{x},(S_{11}-1)e_{y})_{\ell^{2}(\mfrak{X})}=-\overline{(e_{y},JS_{22}Je_{x})_{\ell^{2}(\mfrak{X})}}=-(e_{y},S_{22}e_{x})_{\ell^{2}(\mfrak{X})}
\end{gather*}
for any $x,y\in\mfrak{X}$, which implies the anti-symmetry (\ref{eq:anti-sym_matrix_func}).
\end{Remark}

\setcounter{Theorem}{7}\setcounter{equation}{4}
% Original source lines 267--284
An interesting subclass of PfPPs obtained in this manner consists of those associated with projection operators.
We write the collection of projection operators in $\mcal{Q}(\mcal{K},\Gamma)$ as
\begin{equation*}
	\mrm{Gr}(\mcal{K},\Gamma)=\big\{P\in\mcal{Q}(\mcal{K},\Gamma)\,|\,P^{2}=P\big\}.
\end{equation*}
This notation is, of course, motivated by the fact that a projection operator $P\in\mrm{Gr}(\mcal{K},\Gamma)$ determines a closed subspace $P\mcal{K}\subset\mcal{K}$ and, therefore, the collection of projection operators can be regarded as an analogue of the Grassmann variety.
Let $P_{0}\in \mrm{Gr}(\mcal{K},\Gamma)$ be the projection operator onto the first component of the direct sum decomposition $\mcal{K}=\ell^{2}(\mfrak{X})\oplus\ell^{2}(\mfrak{X})$, which is expressed as
\begin{equation}
\label{eq:projection_vacuum}
	P_{0}=\left(
	\begin{matrix}
		I & 0 \\
		0 & 0
	\end{matrix}
	\right).
\end{equation}
For each $n\in\mbb{Z}_{\ge 0}$, we write $\Omega_{n}=\{\omega\in \Omega\,|\, \#\omega=n\}$ for the collection of $n$-point configurations and set $\Omega^{\circ}=\bigcup_{n=0}^{\infty}\Omega_{n}$, which consists of configurations of finitely many points.


\setcounter{subsection}{2}\setcounter{Theorem}{8}\setcounter{equation}{5}
% Original source lines 289--380
\subsection{CAR algebra and quasifree states}
We introduce an algebra of canonical anti-commutation relations (CAR algebra, for short) with a general one-particle Hilbert space following \cite{Araki1970,Binnenhei1995}.
Another style of, but equivalent, definition of a CAR algebra, can be found in, e.g., \cite{BratteliRobinson1997} (see Remark~\ref{rem:equivalence_CAR} below for the equivalence).
Let~$\mcal{K}$ be a complex Hilbert space of infinite dimension and $\Gamma$ be an anti-unitary involution on $\mcal{K}$.
The algebra $\mcal{C}_{0}(\mcal{K},\Gamma)$ is a~$\ast$-algebra over $\mbb{C}$ generated by $B(f)$, $f\in\mcal{K}$ subject to the relations
\begin{alignat*}{3}
	&B(\alpha f+\beta g)=\alpha B(f)+\beta B(g),\qquad && f, g\in\mcal{K},\quad \alpha,\beta\in\mbb{C},& \\
	&B(f)^{\ast}= B(\Gamma f),\qquad && f\in\mcal{K},&\\
	&\{B(f)^{\ast}, B(g)\}=(f,g)_{\mcal{K}}, \qquad && f,g\in\mcal{K},&
\end{alignat*}
where $\{\cdot,\cdot\}$ is the anti-commutator; $\{a,b\}:=ab+ba$.
It is known that the algebra $\mcal{C}_{0}(\mcal{K},\Gamma)$ admits a unique C$^{\ast}$-norm $\|\cdot\|$. We denote the C$^{\ast}$-completion by $\mcal{C}(\mcal{K},\Gamma)=\overline{\mcal{C}_{0}(\mcal{K},\Gamma)}^{\|\cdot\|}$ and call it the (self-dual) CAR algebra with one-particle space $(\mcal{K},\Gamma)$. To those who are more familiar with regarding $\ell^{2}(\mfrak{X})$ as a one-particle Hilbert space, we emphasize that we adopt a~``doubled'' space as a one-particle Hilbert space.

For a general C$^{\ast}$-algebra $\mfrak{A}$, a state over it is, by definition, a linear functional $\varphi\colon \mfrak{A}\to\mbb{C}$ satisfying the conditions that
\begin{enumerate}\itemsep=0pt
\item[1)] 	for every $A\in\mfrak{A}$, $\varphi (A^{\ast}A)\ge 0$ holds,
\item[2)] 	it is normalized:
		\begin{equation*}
			\|\varphi\|:=\sup\big\{|\varphi (A)|\, \big|\, \|A\|=1\big\}=1.
		\end{equation*}
\end{enumerate}
Note that, from these properties, it can be deduced that $\varphi (1)=1$ (see, e.g., \cite[Chapter~I, Section~9]{Takesaki1979}). Since any positive element $B\in\mfrak{A}$ admits an expression $B=A^{\ast}A$ with some $A\in\mfrak{A}$, the first condition is equivalently stated that $\varphi (B)\ge 0$ for all positive elements $B\in \mfrak{A}$.

\begin{Definition}A state $\varphi$ over a CAR algebra $\mcal{C}(\mcal{K},\Gamma)$ is said to be quasi-free if the moments admit Wick's formula, i.e.,
\begin{gather*}%\label{eq:Wick-property1}
	\varphi (B(f_{1})\cdots B(f_{2n+1}) ) =0, \\
%\label{eq:Wick-property2}
	\varphi (B(f_{1})\cdots B(f_{2n}) ) =(-1)^{n(n-1)/2}\sum_{\sigma}\operatorname{sgn} \sigma \prod_{i=1}^{n}\varphi (B(f_{\sigma (i)})B(f_{\sigma (i+n)}) ),
\end{gather*}
where $\sigma$ runs over permutations in $\mfrak{S}_{2n}$ such that $\sigma(1)<\cdots<\sigma (n)$ and $\sigma (i)<\sigma (i+n)$, $i=1,\dots, n$.
\end{Definition}

It is immediate from the definition that, for a quasi-free state $\varphi$ over $\mcal{C}(\mcal{K},\Gamma)$, a $2n$-point correlation function is expressed in terms of a Pfaffian so that
\begin{equation*}
	\varphi (B(f_{1})\cdots B(f_{2n}) )=\pf \mbb{A}_{\varphi}(f_{1},\dots, f_{2n}),
\end{equation*}
where $\mbb{A}_{\varphi}(f_{1},\dots, f_{2n})$ is the unique anti-symmetric $(2n\times 2n)$-matrix defined by
\begin{equation*}
	\mbb{A}_{\varphi}(f_{1},\dots, f_{2n})_{i,j}=\varphi (B(f_{i})B(f_{j}) ),\qquad 1\le i <j \le 2n.
\end{equation*}

By definition, a quasi-free state over a CAR algebra is uniquely determined by the two-point function. In fact, we have the following.
\begin{Lemma}[\cite{Araki1970}]\label{lem:quasi-free_covariance}
The collection of quasi-free states over $\mcal{C}(\mcal{K},\Gamma)$ is in one-to-one correspondence with the collection of operators $\mcal{Q}(\mcal{K},\Gamma)$ defined in~\eqref{eq:collection_covariance_ops}, under which an operator $S\in\mcal{Q}(\mcal{K},\Gamma)$ corresponds to the quasi-free state~$\varphi_{S}$ defined by
\begin{equation*}%\label{eq:quasifree_covariance}
	\varphi_{S}\big(B(f)^{\ast}B(g)\big)=(f,Sg)_{\mcal{K}},\qquad f,g\in\mcal{K}.
\end{equation*}
\end{Lemma}
We will give a proof of Lemma~\ref{lem:quasi-free_covariance} in Section~\ref{sect:existence_PfPP} for readers' convenience.

As we have announced, we work on the case when $\mcal{K}=\ell^{2}(\mfrak{X})\oplus\ell^{2}(\mfrak{X})$ equipped with $\Gamma$ given by (\ref{eq:Gamma}), and will adopt this pair $(\mcal{K},\Gamma)$ in the sequel without any specification.
The associated CAR algebra $\mcal{C}(\mcal{K},\Gamma)$ is generated by $a_{x}=B((0,e_{x}))$ and $a_{x}^{\ast}=B((e_{x},0))$, $x\in\mfrak{X}$.
Notice that the notation is compatible with the $\ast$-involution since $\Gamma (e_{x},0)=(0,e_{x})$.
Then, from the definition of a quasi-free state, we have
\begin{equation}
\label{eq:quasifree_correlation_function}
	\varphi_{S}\big(a_{x_{1}}^{\ast} \cdots a_{x_{n}}^{\ast}a_{x_{n}}\cdots a_{x_{1}}\big)=\pf [\mbb{K}_{S}(x_{i},x_{j}) ]_{1\le i,j\le n},
\end{equation}
where the matrix-valued function $\mbb{K}_{S}(\cdot,\cdot)\colon \mfrak{X}\times\mfrak{X}\to M(2;\mbb{C})$ was defined in~(\ref{eq:matrix_kernel_S}) associated with the operator~$S$. In fact, the values of the matrix-valued function $\mbb{K}_{S}(\cdot,\cdot)$ is written in terms of the quasi-free state as
\begin{equation*}
	\mbb{K}_{S}(x,y)=\left(
	\begin{matrix}
	\varphi_{S}(a_{x}^{\ast}a_{y}^{\ast}) & \varphi_{S}(a_{x}^{\ast}a_{y}) \\
	\varphi_{S}(a_{x}a_{y}^{\ast})-\delta_{x,y} & \varphi_{S}(a_{x}a_{y})
	\end{matrix}
	\right),\qquad x,y\in\mfrak{X}.
\end{equation*}
At this stage, we can find an analogy between this expectation value and a correlation function of a PfPP.

\begin{Remark}\label{rem:equivalence_CAR}
We may define functionals $a^{\ast}\colon \ell^{2}(\mfrak{X})\to \mcal{C}(\mcal{K},\Gamma)$ and $a\colon \ell^{2}(\mfrak{X})\to \mcal{C}(\mcal{K},\Gamma)$ by extending the assignments $e_{x}\mapsto a_{x}^{\ast}$ and $e_{x}\mapsto a_{x}$, $x\in\mfrak{X}$ linearly and anti-linearly, respectively. In other words, these functionals are defined by $a^{\ast}(u)=B((u,0))$, $a(u)=B((0,Ju))$, $u\in \ell^{2}(\mfrak{X})$.
Then, we find anti-commutation relations that are widely accepted~\cite{BratteliRobinson1997}:
\begin{gather}\label{eq:anti-commutation_BratteliRobinson}
	 \{a^{\ast}(u),a^{\ast}(v)\}=\{a(u),a(v)\}=0, \qquad
	 \{a(u),a^{\ast}(v)\}=(u,v)_{\ell^{2}(\mfrak{X})}, \qquad u,v\in \ell^{2}(\mfrak{X}).
\end{gather}
Conversely, when we have linear and anti-linear functional $a^{\ast}$, $a$ satisfying the anti-commutation relations (\ref{eq:anti-commutation_BratteliRobinson}), the generators of $\mcal{C}(\mcal{K},\Gamma)$ are recovered as $B((u,v))=a^{\ast}(u)+a(Jv)$, $(u,v)\in\mcal{K}$.
\end{Remark}

Let us consider a commutative algebra topologically generated by $a_{x}^{\ast}a_{x}$, $x\in\mfrak{X}$, which is identified with the algebra $C(\Omega)$ of continuous functions on $\Omega$ by the correspondence
\begin{equation*}
	\prod_{i=1}^{n}a_{x_{i}}^{\ast}a_{x_{i}}\mapsto \chi_{\Omega_{x_{1},\dots, x_{n}}},\qquad x_{1},\dots, x_{n}\in\mfrak{X}\colon \ \mbox{distinct.}
\end{equation*}
In the sequel, we regard $C(\Omega)$ as a subalgebra of $\mcal{C}(\mcal{K},\Gamma)$ under this correspondence.

Our strategy to prove Proposition~\ref{prop:existence_ppp} is to identify the above expectation value~(\ref{eq:quasifree_correlation_function}) with the correlation function of the desired PfPP, expecting that for a quasi-free state $\varphi_{S}$, there exists a probability measure $M^{S}$ on $(\Omega,\Sigma)$ so that the restriction of $\varphi_{S}$ on the subalgebra $C(\Omega)$ is identical to the integration with respect to $M^{S}$:
\begin{equation*}
	\varphi_{S}(F)=\int_{\Omega}F(\omega)M^{S}({\rm d}\omega),\qquad F\in C(\Omega).
\end{equation*}
Then, the probability measure $M^{S}$ is automatically a PfPP.
We will see in Section~\ref{sect:existence_PfPP} that this indeed happens to prove Proposition \ref{prop:existence_ppp}.


\setcounter{subsection}{5}\setcounter{Theorem}{14}\setcounter{equation}{7}
% Original source lines 442--501
\subsection{Conditional measures of PfPPs}
Let $X$ and $\pr{X}$ be disjoint finite subsets in $\mfrak{X}$ and take a cylinder set
\begin{equation*}
	C(X,\pr{X})=\{\omega\in \Omega\,|\, X\subset \omega,\, \pr{X}\cap\omega=\varnothing\},
\end{equation*}
which consists of configurations such that every points in $X$ are occupied and those in $\pr{X}$ are unoccupied.
We identify the cylinder set $C(X,\pr{X})$ with $\Omega (\mfrak{X}\backslash (X\sqcup\pr{X})):=\{0,1\}^{\mfrak{X}\backslash (X\sqcup\pr{X})}$ via the map
\begin{equation*}
	F_{X,\pr{X}}\colon \ C(X,\pr{X})\to \Omega (\mfrak{X}\backslash (X\sqcup\pr{X}));\qquad \omega\mapsto \omega\backslash X.
\end{equation*}
For a probability measure $M$ on $(\Omega,\Sigma)$, assume that the cylinder set has strictly positive weight: $M(C(X,\pr{X}))>0$.
We define the conditional measure of $M$ on $(X,\pr{X})$ by
\begin{equation*}
	M_{X,\pr{X}}:=(F_{X,\pr{X}})_{\ast}\frac{M\big|_{C(X,\pr{X})}}{M(C(X,\pr{X}))}.
\end{equation*}

We focus on PfPPs associated with projection operators.
For a subset $A\subset\mfrak{X}$, we set
\begin{equation*}
	\mcal{K}_{A}^{+}=\overline{\mrm{Span}\{(e_{x},0)\,|\,x\in A\}},\qquad \mcal{K}^{-}_{A}=\overline{\mrm{Span}\{(0,e_{x})\,|\,x\in A\}}
\end{equation*}
and $\mcal{K}_{A}=\mcal{K}_{A}^{+}\oplus\mcal{K}_{A}^{-}$, which is the direct sum of Hilbert spaces.
It is obvious from the definition that $\Gamma$ preserves the subspace $\mcal{K}_{A}$. Thus, we can write $\Gamma_{A}:=\Gamma|_{\mcal{K}_{A}}$.

Let $X$ and $\pr{X}$ be finite disjoint subsets of $\mfrak{X}$ and take a projection operator $P\in \mrm{Gr}(\mcal{K},\Gamma)$.
Then, $P\mcal{K}\subset\mcal{K}$ is a closed subspace. We define a new closed subspace
\begin{equation*}
	P\mcal{K}_{X,\pr{X}}:=\big(P\mcal{K}+\big(\mcal{K}_{X}^{+}\oplus \mcal{K}^{-}_{\pr{X}}\big)\big)\cap \mcal{K}_{\mfrak{X}\backslash (X\sqcup\pr{X})}
\end{equation*}
in $\mcal{K}_{\mfrak{X}\backslash (X\sqcup\pr{X})}$ and a projection operator $P_{X,\pr{X}}$ as the orthogonal projection onto $P\mcal{K}_{X,\pr{X}}$ in $\mcal{K}_{\mfrak{X}\backslash (X\sqcup\pr{X})}$.

\begin{Lemma}\label{lem:conditional_projection}
Let $X$ and $\pr{X}$ be finite disjoint subsets of $\mfrak{X}$. Then, we have
\begin{equation*}
	P_{X,\pr{X}}\in\mrm{Gr}\big(\mcal{K}_{\mfrak{X}\backslash(X\sqcup\pr{X})},\Gamma_{\mfrak{X}\backslash (X\sqcup\pr{X})}\big).
\end{equation*}
\end{Lemma}

 Let us also introduce the notion of regularity.
 \begin{Definition} Let $X$ and $\pr{X}$ be finite disjoint subsets of $\mfrak{X}$.
 A projection operator $P\in \mrm{Gr}(\mcal{K},\Gamma)$ is said to be $(X,\pr{X})$-regular if
 \begin{equation*}
 	P\mcal{K}\cap \big(\mcal{K}_{X}^{+}\oplus \mcal{K}_{\pr{X}}^{-}\big)=\{0\}.
 \end{equation*}
 \end{Definition}

 The following result is an analogue of \cite[Proposition~6.13]{Olshanski2020} for PfPPs.
 \begin{Theorem}\label{thm:conditional_projection}
 Let $X$ and $\pr{X}$ be finite disjoint subsets of $\mfrak{X}$.
 Assume that a projection operator $P\in\mrm{Gr}(\mcal{K},\Gamma)$ is $(X,\pr{X})$-regular.
 Then, we have
 \begin{equation*}
 	M^{P}_{X,\pr{X}}:=\big(M^{P}\big)_{X,\pr{X}}=M^{P_{X,\pr{X}}}.
 \end{equation*}
 In particular, it is a PfPP associated with a projection operator.
 \end{Theorem}

 \begin{Remark}
A similar problem has been studied in~\cite{BufetovCundenQiu2019}. Our result describes the reduction of a projection operator, which is new.
 \end{Remark}

\setcounter{section}{1}
% Original source lines 511--686
\section{Existence of Pfaffian point processes} \label{sect:existence_PfPP}
\subsection{Fock representations}
Here, we see Fock representations of $\mcal{C}(\mcal{K},\Gamma)$, which play prominent roles in the theory of a CAR algebra.
\subsubsection{General construction}
Let $\mcal{H}$ be a Hilbert space.
For each $n\in\mbb{N}$, we write $\bigwedge^{n}\mcal{H}$ for the $n$-th wedge product of $\mcal{H}$, which is generated by the vectors $u_{1}\wedge\cdots \wedge u_{n}$, $u_{1},\dots, u_{n}\in\mcal{H}$ subject to the anti-symmetry:
\begin{equation*}
	u_{1}\wedge\cdots \wedge u_{n}=\operatorname{sgn} (\sigma) u_{\sigma(1)}\wedge\cdots \wedge u_{\sigma (n)},\qquad \sigma\in \mfrak{S}_{n}.
\end{equation*}
When we take $\{e_{i}\}_{i=1,2,\dots}$ for a complete orthonormal system of $\mcal{H}$, then the vectors $e_{i_{1}}\wedge\cdots \wedge e_{i_{n}}$, $i_{1}>\cdots >i_{n}$ form a complete orthonormal system of $\bigwedge^{n}\mcal{H}$.
The Fermi Fock space over $\mcal{H}$ is defined by
\begin{equation*}
	\mcal{F}(\mcal{H})=\bigoplus_{n=0}^{\infty}\bigwedge^{n}\mcal{H},
\end{equation*}
where we set $\bigwedge^{0}\mcal{H}=\mbb{C}\bm{1}$.
The Fermi Fock space admits the natural inner product induced from each component of the direct sum and becomes a Hilbert space, i.e., the direct sum is understood in the topological sense.

For each $u\in \mcal{H}$, the creation operator $a_{\mcal{H}}^{\ast}(u)$ is an operator on $\mcal{F}(\mcal{H})$ defined by
\begin{equation*}
	a_{\mcal{H}}^{\ast}(u)\colon \ \eta\mapsto u\wedge \eta,\qquad \eta\in \mcal{F}(\mcal{H}).
\end{equation*}
The annihilation operator $a_{\mcal{H}}(u)$ is defined as the adjoint operator of $a_{\mcal{H}}^{\ast}(u)$.
By definition, it is obvious that creation and annihilation operators act as
\begin{equation*}
	a_{\mcal{H}}^{\ast}(u)\colon \ \bigwedge^{n}\mcal{H}\to \bigwedge^{n+1}\mcal{H},\qquad a_{\mcal{H}}(u)\colon \ \bigwedge^{n}\mcal{H}\to \bigwedge^{n-1}\mcal{H}.
\end{equation*}
We can also see that the assignment $u\mapsto a_{\mcal{H}}^{\ast}(u)$ is linear and $u\mapsto a_{\mcal{H}}(u)$ is anti-linear.

\subsubsection{Fock representation and a quasi-free state}
Let us take a projection operator $P\in\mrm{Gr}(\mcal{K},\Gamma)$, associated to which we can construct a representation of $\mcal{C}(\mcal{K},\Gamma)$ on a Fock space $\mcal{F}(P\mcal{K})$.
To describe the action, notice that, from the property $1-P=\overline{P}$, we can see that the projection to the complementary subspace of $P\mcal{K}$ is $\overline{P}$.
Let us set
\begin{equation*}
	\pi_{P}\colon \ \mcal{C}(\mcal{K},\Gamma)\to \mfrak{B}(\mcal{F}(P\mcal{K}));\qquad B(f)\mapsto a_{P\mcal{K}}^{\ast}(Pf)+a_{P\mcal{K}}(P\Gamma f),\qquad f\in\mcal{K}.
\end{equation*}
Then, it is known that $(\pi_{P},\mcal{F}(P\mcal{K}))$ is a faithful and irreducible representation of $\mcal{C}(\mcal{K},\Gamma)$.
When we set
\begin{equation*}
	\varphi_{P}(A)=(\bm{1}, \pi_{P}(A)\bm{1})_{\mcal{F}(P\mcal{K})},\qquad A\in \mcal{C}(\mcal{K},\Gamma),
\end{equation*}
we can verify that $\varphi_{P}$ is exactly the quasi-free state corresponding to $P$ in the sense that it possesses the property required in Lemma \ref{lem:quasi-free_covariance}

In particular, when we take $P_{0}\in\mrm{Gr}(\mcal{K},\Gamma)$, we have $P_{0}\mcal{K}\simeq \ell^{2}(\mfrak{X})$, and the map $\pi_{P_{0}}$ is described as
\begin{equation*}
	\pi_{P_{0}}(a^{\ast}(u))=a_{P_{0}\mcal{K}}^{\ast}(u),\qquad \pi_{P_{0}}(a(u))=a_{P_{0}\mcal{K}}(u),\qquad u\in P_{0}\mcal{K}\simeq \ell^{2}(\mfrak{X}).
\end{equation*}

Let us describe a standard complete orthonormal system of the Fock space $\mcal{F}\big(\ell^{2}(\mfrak{X})\big)$. Since~$\mfrak{X}$ is countable, it can be equipped with a linear order~$\le$. For each $\omega=\{x_{1}>\cdots>x_{n}\}\in\Omega^{\circ}$, we set
\begin{equation*}
	e_{\omega}:=e_{x_{1}}\wedge\cdots \wedge e_{x_{n}}.
\end{equation*}
Then the collection $\{e_{\omega}\}_{\omega\in\Omega^{\circ}}$ forms a complete orthonormal system.

\begin{proof}[Proof of Lemma \ref{lem:quasi-free_covariance}]
Let $\varphi$ be a quasi-free state over $\mcal{C}(\mcal{K},\Gamma)$.
Then the assignment
\begin{equation*}
	Q_{\varphi}\colon \ \mcal{K}\times \mcal{K}\to\mbb{C};\qquad (f,g)\mapsto \varphi (B(f)^{\ast}B(g))
\end{equation*}
defines a quadratic form on $\mcal{K}$.
It follows from the relation in the CAR algebra that
\begin{equation*}
	B(f)^{\ast}B(f)\le B(f)^{\ast}B(f)+B(f)B(f)^{\ast}=\|f\|^{2},\qquad f\in\mcal{K}.
\end{equation*}
Since the norm of a state is unity, we have $Q_{\varphi}(f,f)\le \|B(f)^{\ast}B(f)\|\le \|f\|^{2}$ for any $f\in\mcal{K}$, which implies that $Q_{\varphi}$ is a bounded quadratic form. Therefore, owing to the correspondence between bounded operators and bounded quadratic forms (see, e.g., \cite[Chapter~1]{Koshmanenko1999}), there exists a bounded operator $S\in \mfrak{B}(\mcal{K})$ such that $Q_{\varphi}(f,g)=(f,Sg)_{\mcal{K}}$, $f,g \in \mcal{K}$. It also follows from the positivity of the state that the quadratic form $Q_{\varphi}$ is positive, and therefore, $S=S^{\ast}\ge 0$. Again, from the relation in the CAR algebra, we have
\begin{equation*}
	(f,Sg)_{\mcal{K}}=(f,g)_{\mcal{K}}-(\Gamma g, S\Gamma f)_{\mcal{K}}=(f,g)_{\mcal{K}}-\big(f,\overline{S} g\big)_{\mcal{K}},\qquad f,g\in\mcal{K},
\end{equation*}
which implies that $\Gamma S \Gamma=1-S$. Since $\Gamma S \Gamma\ge 0$, we can see that $S\le 1$.

Conversely, given an operator $S\in\mcal{Q}(\mcal{K},\Gamma)$, we define the following operator
\begin{equation*}
	P_{S}=\left(
	\begin{matrix}
	S & S^{1/2}(1-S)^{1/2} \\
	S^{1/2}(1-S)^{1/2} & 1-S
	\end{matrix}
	\right)
\end{equation*}
acting on $\what{\mcal{K}}=\mcal{K}\oplus\mcal{K}$. Notice that the assumption $0\le S=S^{\ast}\le 1$ ensures that the square roots~$S^{1/2}$ and $(1-S)^{1/2}$ make sense. Let us equip this Hilbert space with the anti-unitary involution
\begin{equation*}
	\what{\Gamma}=\left(
	\begin{matrix}
	\Gamma & 0 \\
	0 & -\Gamma
	\end{matrix}
	\right).
\end{equation*}
Then, it can be checked that $P_{S}\in \mrm{Gr}\big(\what{\mcal{K}},\what{\Gamma}\big)$, which implies that the functional $\varphi_{P_{S}}$ defined by
\begin{equation*}
	\varphi_{P_{S}}(A)= (\bm{1},\pi_{P_{S}}(A)\bm{1} )_{\mcal{F}(P_{S}\what{\mcal{K}})},\qquad A\in \mcal{C}\big(\what{\mcal{K}},\what{\Gamma}\big)
\end{equation*}
is a quasi-free state over $\mcal{C}\big(\what{\mcal{K}},\what{\Gamma}\big)$. Now, since $\what{\Gamma}$ acts diagonally along the direct sum decomposition, $\mcal{C}(\mcal{K},\Gamma)$ is regarded as a subalgebra of $\mcal{C}\big(\what{\mcal{K}},\what{\Gamma}\big)$ and the restriction $\varphi_{S}=\varphi_{P_{S}}|_{\mcal{C}(\mcal{K},\Gamma)}$ is the quasi-free state corresponding to the given~$S$.
\end{proof}

\subsection{Proof of Proposition \ref{prop:existence_ppp}}
Now, we are at the position of proving Proposition~\ref{prop:existence_ppp}.
Our strategy is to check the criteria by Lenard \cite{Lenard1975a, Lenard1975b}.

\subsubsection{Lenard's criteria}
Suppose that a system of functions $\big\{\rho_{n}\colon\mfrak{X}^{n}\to\mbb{R}\big\}_{n=1}^{\infty}$ is given.
A question is if there exists a~probability measure~$M$ on $(\Omega,\Sigma)$ such that its $n$-point correlation function $\rho^{M}_{n}$ coincides with the given function $\rho_{n}$ for every $n\in\mbb{N}$.
Lenard \cite{Lenard1975a, Lenard1975b} clarified the necessary and sufficient conditions for this to happen.

\begin{Theorem}[\cite{Lenard1975a, Lenard1975b}]
A system of functions $\big\{\rho_{n}\colon \mfrak{X}^{n}\to\mbb{R}\big\}_{n=1}^{\infty}$ is a one of correlation functions for a probability measure on $(\Omega,\Sigma)$ if it possesses the following properties:
\begin{enumerate}\itemsep=0pt
\item[$1.$] {\it Symmetry:} each function $\rho_{n}$ is a symmetric function, i.e.,
	\begin{equation*}
		\rho_{n}(x_{\sigma (1)},\dots, x_{\sigma (n)})=\rho_{n}(x_{1},\dots, x_{n})
	\end{equation*}
	for arbitrary $\sigma\in\mfrak{S}_{n}$ and $x_{1},\dots, x_{n}\in\mfrak{X}$.
\item[$2.$] {\it Positivity:} for any system of functions $\bm{\Phi}=\big\{\Phi_{n}\colon \mfrak{X}^{n}\to\mbb{R}\big\}_{n=0}^{N}$ $(\Phi_{0}$ is understood as a~constant$)$ with finite support such that
	\begin{equation}\label{eq:positivity_test_function}
		\Phi_{0}+\sum_{n=1}^{N}\sum_{\substack{i_{1},\dots,i_{n}\in I \\ \mrm{distinct}}}\Phi_{n}(x_{i_{1}},\dots, x_{i_{n}})\ge 0
	\end{equation}
	holds for any $\omega=\{x_{i}\}_{i\in I}\in\Omega$, we have
	\begin{equation*}
		\Phi_{0}+\sum_{n=1}^{N}\sum_{x_{1},\dots, x_{n}\in\mfrak{X}}\Phi_{n}(x_{1},\dots, x_{n})\rho_{n}(x_{1},\dots, x_{n})\ge 0.
	\end{equation*}
\end{enumerate}
\end{Theorem}

Therefore, Proposition \ref{prop:existence_ppp} reduces to the following assertion.
\begin{Lemma}
\label{lem:existence_ppp}
Let $S\in\mcal{Q}(\mcal{K},\Gamma)$. Then the system of functions
\begin{equation*}
	\rho_{n}(x_{1},\dots, x_{n})=\pf \left[\mbb{K}_{S}(x_{i},x_{j})\right]_{1\le i,j\le n}, \qquad x_{1},\dots, x_{n}\in\mfrak{X}, \qquad n\in\mbb{N}
\end{equation*}
satisfies the symmetry and positivity conditions.
\end{Lemma}

\subsubsection{Proof of Lemma \ref{lem:existence_ppp}}
As we have seen, each correlation function is realized as
\begin{equation*}
	\rho_{n}(x_{1},\dots, x_{n})=\varphi_{S}\big(a_{x_{1}}^{\ast}\cdots a_{x_{n}}^{\ast}a_{x_{n}}\cdots a_{x_{1}}\big).
\end{equation*}
Then, the symmetry condition is obviously satisfied since
\begin{equation*}
	a_{x_{1}}^{\ast}\cdots a_{x_{n}}^{\ast}a_{x_{n}}\cdots a_{x_{1}}
	=a_{x_{\sigma (1)}}^{\ast}\cdots a_{x_{\sigma (n)}}^{\ast}a_{x_{\sigma (n)}}\cdots a_{x_{\sigma (1)}}
\end{equation*}
holds for all $\sigma\in\mfrak{S}_{n}$.

To show the positivity condition, let $\bm{\Phi}=\big\{\Phi_{n}\colon\mfrak{X}^{n}\to\mbb{R}\big\}_{n=0}^{N}$ be a system of functions satisfying~(\ref{eq:positivity_test_function}).
We can see that
\begin{equation*}
	\Phi_{0}+\sum_{n=1}^{N}\sum_{x_{1},\dots, x_{n}\in\mfrak{X}}\Phi_{n}(x_{1},\dots, x_{n})\rho_{n}(x_{1},\dots, x_{n})
	=\varphi_{S} (A_{\bm{\Phi}} ),
\end{equation*}
where we set
\begin{equation*}
	A_{\bm{\Phi}}=\Phi_{0}+\sum_{n=1}^{N}\sum_{x_{1},\dots,x_{n}\in\mfrak{X}}\Phi_{n}(x_{1},\dots,x_{n})a_{x_{1}}^{\ast}\cdots a_{x_{n}}^{\ast}a_{x_{n}}\cdots a_{x_{1}}.
\end{equation*}
Therefore, owing to the positivity of a state over a C$^{\ast}$-algebra, it suffices to show that $A_{\bm{\Phi}}\in\mcal{C}(\mcal{K},\Gamma)$ is a positive element, which can be checked in any faithful representation. In fact, when we take a faithful representation $(\pi,\mcal{H})$, we may regard $\mcal{C}(\mcal{K},\Gamma)$ as a C$^{\ast}$-subalgebra of $\mfrak{B}(\mcal{H})$ via the embedding~$\pi$. Since the spectrum of an element in a C$^{\ast}$-subalgebra coincides with that in a~whole algebra (see, e.g., \cite[Proposition~2.2.7]{BratteliRobinson1987}), the relevant element $A_{\bm{\Phi}}\in\mcal{C}(\mcal{K},\Gamma)$ is positive if and only if $\pi \left(A_{\bm{\Phi}}\right)$ is a positive operator on $\mcal{H}$.

We can, in particular, take a Fock representation $\big(\pi_{P_{0}},\mcal{F}\big(\ell^{2}(\mfrak{X})\big)\big)$.
\begin{Lemma}
In the Fock space $\mcal{F}\big(\ell^{2}(\mfrak{X})\big)$, the complete orthonormal system $\{e_{\omega}\}_{\omega\in\Omega^{\circ}}$ diagonalizes $\pi_{P_{0}} (A_{\bm{\Phi}} )$ so that
\begin{equation*}
	\pi_{P_{0}}\left(A_{\bm{\Phi}}\right) e_{\omega}=\left(\Phi_{0}+\sum_{k=1}^{N}\sum_{\substack{i_{1},\dots,i_{k}=1\\ \mrm{distinct}}}^{n}\Phi_{k}(x_{i_{1}},\dots, x_{i_{k}})\right) e_{\omega},\qquad \omega=\{x_{1}>\cdots >x_{n}\},
\end{equation*}
\end{Lemma}
\begin{proof}
This follows from a direct computation.
Let us notice that $a_{x_{1}}^{\ast}\cdots a_{x_{n}}^{\ast}a_{x_{n}}\cdots a_{x_{1}}=0$ if any two points from $x_{1},\dots, x_{n}$ coincide. Hence, we have
\begin{equation*}
	A_{\bm{\Phi}}=\Phi_{0}+\sum_{n=1}^{N}\sum_{\substack{x_{1},\dots,x_{n}\in\mfrak{X}\\ \mrm{distinct}}}\Phi_{n}(x_{1},\dots,x_{n})\prod_{i=1}^{n}a_{x_{i}}^{\ast}a_{x_{i}}.
\end{equation*}
We can also see that $\pi_{P_{0}}(a_{x}^{\ast}a_{x})e_{\omega}=\chi_{[x\in\omega]}e_{\omega}$, $x\in\mfrak{X}$, $\omega\in\Omega^{\circ}$. Therefore, the desired result is obtained.
\end{proof}

The eigenvalues of $\pi_{P_{0}} (A_{\bm{\Phi}} )$ are non-negative from the assumption implying that $A_{\bm{\Phi}}$ is a positive element, and therefore, the system of functions $\{\rho_{n}\}_{n=1}^{\infty}$ fulfills the positivity conditions.
Now, the proof is complete.


\setcounter{section}{3}
% Original source lines 1293--1654
\section{Conditional measures}\label{sect:conditional_measures}
\subsection{Proof of Lemma \ref{lem:conditional_projection}}
First, we notice that the subspace $P\mcal{K}_{X,\pr{X}}$ consists of vectors $v\in \mcal{K}_{\mfrak{X}\backslash(X\sqcup\pr{X})}$ such that there exist $a\in \mcal{K}_{X}^{+}$, $\pr{a}\in \mcal{K}_{\pr{X}}^{-}$ and $v+a+\pr{a}\in P\mcal{K}$. Due to the property $\Gamma \mcal{K}_{X}^{\pm}=\mcal{K}_{X}^{\mp}$, we can see that
\begin{equation*}
	\Gamma_{\mfrak{X}\backslash(X\sqcup\pr{X})} P\mcal{K}_{X,\pr{X}}=\big(\overline{P}\mcal{K}+\big(\mcal{K}_{X}^{-}\oplus \mcal{K}_{\pr{X}}^{+}\big)\big)\cap \mcal{K}_{\mfrak{X}\backslash (X\cup \pr{X})}=\overline{P}\mcal{K}_{\pr{X},X}.
\end{equation*}
We can also see that $ (P\mcal{K}_{X,\pr{X}})^{\perp}=\overline{P}\mcal{K}_{\pr{X},X}$. In fact, for $u\in P\mcal{K}_{X,\pr{X}}$ and $v\in \overline{P}\mcal{K}_{\pr{X},X}$, we can take $a\in \mcal{K}_{X}^{+}$, $\pr{a}\in \mcal{K}_{\pr{X}}^{-}$, $b\in \mcal{K}_{X}^{-}$, and $\pr{b}\in \mcal{K}_{\pr{X}}^{+}$ such that
\begin{gather*}
	u+a+\pr{a} \in P\mcal{K}, \qquad v+b+\pr{b}\in \overline{P}\mcal{K}.
\end{gather*}
Now, since $a$, $\pr{a}$, $b$, $\pr{b}$ are mutually orthogonal and orthogonal to $u$ and $v$, we have
\begin{equation*}
	(u,v)_{\mcal{K}_{\mfrak{X}\backslash (X\cup\pr{X})}}=(u+a+\pr{a},v+b+\pr{b})_{\mcal{K}}=0,
\end{equation*}
which implies $(P\mcal{K}_{X,\pr{X}})^{\perp}=\overline{P}\mcal{K}_{\pr{X},X}$.
Therefore, $\overline{P}_{X,\pr{X}}=1-P_{X,\pr{X}}$ holds.

\subsection{Conditioning on the CAR algebra}\label{subsect:conditioning_CAR_alg}
For finite disjoint subsets $X, \pr{X}\subset\mfrak{X}$, the characteristic function on the corresponding cylinder set $C(X,\pr{X})$ is identified as
\begin{equation*}
	\chi_{C(X,\pr{X})}=\prod_{x\in X}a_{x}^{\ast}a_{x} \prod_{x\in \pr{X}} a_{x}a_{x}^{\ast}\in C(\Omega)\subset \mcal{C}(\mcal{K},\Gamma).
\end{equation*}
In general, given a state $\varphi$ over $\mcal{C}(\mcal{K},\Gamma)$ its conditioning over $C(X,\pr{X})$ is a state defined by (see, e.g.,~\cite{RedeiSummers2006})
\begin{equation*}
	\varphi (A|\chi_{C(X,\pr{X})})=\frac{\varphi (\chi_{C(X,\pr{X})}A\chi_{C(X,\pr{X})})}{\varphi (\chi_{C(X,\pr{X})})},\qquad A\in \mcal{C}(\mcal{K},\Gamma).
\end{equation*}
According to the decomposition $\mfrak{X}=X\sqcup \pr{X}\sqcup \mfrak{X}\backslash (X\sqcup\pr{X})$, we have a decomposition
\begin{equation*}
	\mcal{C}(\mcal{K},\Gamma)=\mcal{C}(\mcal{K}_{X\sqcup \pr{X}},\Gamma_{X\sqcup\pr{X}})\otimes \mcal{C}(\mcal{K}_{\mfrak{X}\backslash (X\sqcup\pr{X})}),
\end{equation*}
which enables us to regard $C(\mcal{K}_{\mfrak{X}\backslash (X\sqcup\pr{X})},\Gamma_{\mfrak{X}\backslash (X\sqcup\pr{X})})$ as a subalgebra of $\mcal{C}(\mcal{K},\Gamma)$ that coincides with the subalgebra realized as $\chi_{C(X,\pr{X})}\mcal{C}(\mcal{K},\Gamma)\chi_{C(X,\pr{X})}$.
Therefore, the conditional state $\varphi (\cdot|\chi_{C(X,\pr{X})})$ is regarded as a state over $C(\mcal{K}_{\mfrak{X}\backslash (X\sqcup\pr{X})},\Gamma_{\mfrak{X}\backslash (X\sqcup\pr{X})})$ and computed as
\begin{equation*}
	\varphi (A|\chi_{C(X,\pr{X})})=\frac{\varphi (A\chi_{C(X,\pr{X})})}{\varphi (\chi_{C(X,\pr{X})})},\qquad A\in C(\mcal{K}_{\mfrak{X}\backslash (X\sqcup\pr{X})},\Gamma_{\mfrak{X}\backslash (X\sqcup\pr{X})}).
\end{equation*}
In particular, in the case when $\varphi=\varphi_{S}$ is a quasi-free state associated with $S\in \mcal{Q}(\mcal{K},\Gamma)$, we have
\begin{equation*}
	\varphi_{S}(F|\chi_{C(X,\pr{X})})=\int_{\Omega (\mfrak{X}\backslash (X\sqcup\pr{X}))}F(\omega)M^{S}_{X,\pr{X}}({\rm d}\omega),\qquad F\in C(\Omega (\mfrak{X}\backslash (X\sqcup\pr{X}))).
\end{equation*}

\subsection{Some observations}
We see that obtaining $P_{X,\pr{X}}$ from a projection operator $P\in\mrm{Gr}(\mcal{K},\Gamma)$ is decomposed into fundamental steps.
\begin{Lemma}
\label{lem:rest_decomp} We have the following decomposition properties.
\begin{enumerate}\itemsep=0pt
\item[$1.$] 	When $X, \pr{X}\subset\mfrak{X}$ are disjoint finite subsets, we have
		\begin{equation}		\label{eq:rest_proj_decomp}
			P_{X,\pr{X}}=(P_{\varnothing,\pr{X}})_{X,\varnothing}=(P_{X,\varnothing})_{\varnothing, \pr{X}}.
		\end{equation}
\item[$2.$] 	When $X=X_{1}\sqcup X_{2}\subset \mfrak{X}$ is a disjoint union of finite subsets, then
		\begin{equation*}
			P_{X,\varnothing}=(P_{X_{1},\varnothing})_{X_{2},\varnothing}.
		\end{equation*}
\item[$3.$] 	When $\pr{X}=\pr{X}_{1}\sqcup \pr{X}_{2}\subset\mfrak{X}$ is a disjoint union of finite subsets, then
		\begin{equation*}
			P_{\varnothing,\pr{X}}=(P_{\varnothing,\pr{X}_{1}})_{\varnothing,\pr{X}_{2}}.
		\end{equation*}
\end{enumerate}
\end{Lemma}
\begin{proof}We only prove (\ref{eq:rest_proj_decomp}) since the other two follow from arguments of the same type.
We have noticed that $P\mcal{K}$ consists of vectors $u\in \mcal{K}_{\mfrak{X}\backslash (X\sqcup\pr{X})}$ such that there exist $a\in \mcal{K}_{X}^{+}$, $\pr{a}\in \mcal{K}_{\pr{X}}^{-}$ and $u+a+\pr{a}\in P\mcal{K}$. This exactly means that
\begin{equation*}
	u+a\in \{v\in \mcal{K}_{\mfrak{X}\backslash \pr{X}}\,|\, \exists\, \pr{a}\in \mcal{K}_{\pr{X}},\, v+\pr{a}\in P\mcal{K}\},
\end{equation*}
while the latter set is just $P\mcal{K}_{\varnothing,\pr{X}}$. Therefore, $P\mcal{K}_{X,\pr{X}}=(P\mcal{K}_{\varnothing,\pr{X}})_{X,\varnothing}$. The other relation $P\mcal{K}_{X,\pr{X}}=(P\mcal{K}_{X,\varnothing})_{\varnothing,\pr{X}}$ is also shown. Consequently, we reach~(\ref{eq:rest_proj_decomp}).
\end{proof}

We also notice that the regularity is inherited along decomposition.
\begin{Lemma}\label{lem:regularity_decomp} We have the following properties regarding the regularity.
\begin{enumerate}\itemsep=0pt
\item[$1.$] 	For disjoint finite subsets $X,\pr{X}\subset\mfrak{X}$, a projection operator $P\in\mrm{Gr}(\mcal{K},\Gamma)$ is $(X,\pr{X})$-regular if and only if one of the following equivalent conditions holds:
		\begin{enumerate}\itemsep=0pt
		\item[$(a)$] 	$P$ is $(X,\varnothing)$-regular and $P_{X,\varnothing}$ is $(\varnothing,\pr{X})$-regular.
		\item[$(b)$] 	$P$ is $(\varnothing,\pr{X})$-regular and $P_{\varnothing,\pr{X}}$ is $(X,\varnothing)$-regular.
		\end{enumerate}
\item[$2.$] 	When $X=X_{1}\sqcup X_{2}$ is a disjoint union of finite subsets, a projection operator $P\in\mrm{Gr}(\mcal{K},\Gamma)$ is $(X,\varnothing)$-regular if and only if $P$ is $(X_{1},\varnothing)$-regular and $P_{X_{1},\varnothing}$ is $(X_{2},\varnothing)$-regular.
\item[$3.$] 	When $\pr{X}=\pr{X}_{1}\sqcup \pr{X}_{2}$ is a disjoint union of finite subsets, a projection operator $P\in\mrm{Gr}(\mcal{K},\Gamma)$ is $(\varnothing,\pr{X})$-regular if and only if $P$ is $(\varnothing,\pr{X}_{1})$-regular and $P_{\varnothing,\pr{X}_{1}}$ is $(\varnothing,X_{2})$-regular.
\end{enumerate}
\end{Lemma}
\begin{proof}
We only prove the equivalence to the item (a) in (1) since the others are shown in similar arguments.
Suppose that $P\in \mrm{Gr}(\mcal{K},\Gamma)$ is $(X,\pr{X})$-regular. Then, it is obvious that $P$ is $(X,\varnothing)$-regular. Now, $P\mcal{K}_{X,\varnothing}$ consists of vectors $u\in\mcal{K}_{\mfrak{X}\backslash X}$ such that there exists $a\in \mcal{K}_{X}^{+}$ and $u+a\in P\mcal{K}$. Let us take a vector $u_{0}\in P\mcal{K}_{X,\varnothing}\cap \mcal{K}_{\pr{X}}^{-}$. Then, there exists $a_{0}\in \mcal{K}_{X}^{+}$ such that $u_{0}+a_{0}\in P\mcal{K}$, while we also have $u_{0}+a_{0}\in \mcal{K}_{X}^{+}\oplus \mcal{K}_{\pr{X}}^{-}$, which implies that $u_{0}+a_{0}\in P\mcal{K}\cap \mcal{K}_{X}^{+}\oplus \mcal{K}_{\pr{X}}^{-}=\{0\}$ by assumption. Therefore, we must have $u_{0}=0$ and see that $P\mcal{K}_{X,\varnothing}$ is $(\varnothing,\pr{X})$-regular.

Conversely, let us suppose that $P$ is $(X,\varnothing)$-regular and $P_{X,\varnothing}$ is $(\varnothing,\pr{X})$-regular. Take a~nonzero vector $u_{1}\in P\mcal{K}\cap \mcal{K}_{X}^{+}\oplus \mcal{K}_{\pr{X}}^{-}$. Since $P$ is $(X,\varnothing)$-regular, when we write $u_{1}=a_{1}+\pr{a}_{1}$, $a_{1}\in \mcal{K}_{X}^{+}$, $\pr{a}\in \mcal{K}_{\pr{X}}^{-}$, we must have $\pr{a}_{1}\neq 0$. On the other hand, by definition, $\pr{a}_{1}\in P\mcal{K}_{X,\varnothing}$, which implies that $\pr{a}_{1}\in P\mcal{K}_{X,\varnothing}\cap \mcal{K}_{\pr{X}}^{-}$ contradicting the assumption that $P_{X,\varnothing}$ is $(\varnothing,\pr{X})$-regular. Therefore, $P$ is $(X,\pr{X})$-regular.
\end{proof}

Lemmas \ref{lem:rest_decomp} and \ref{lem:regularity_decomp} verify that it suffices to prove Theorem \ref{thm:conditional_projection} in the case when $(X,\pr{X})=(\{x\},\varnothing)$ and $(\varnothing,\{x\})$ with $x\in\mfrak{X}$.
Let us introduce subspaces
\begin{gather*}
	\mcal{R}_{\pm}(x;\mcal{K}):=\big\{u\in \mcal{K}\,|\, \big(\mcal{K}^{\pm}_{\{x\}},u\big)_{\mcal{K}}=0\big\},\qquad x\in \mfrak{X},
\end{gather*}
and set $\mcal{R}_{\pm}(x;P\mcal{K})=P\mcal{K}\cap \mcal{R}_{\pm}(x;\mcal{K})$.
When we write an element $u\in P\mcal{K}$ as $u=(f,g)$ according to the direct sum decomposition $\mcal{K}=\ell^{2}(\mfrak{X})\oplus \ell^{2}(\mfrak{X})$, these are just
\begin{gather*}
	\mcal{R}_{+}(x;P\mcal{K}) =\{(f,g)\in P\mcal{K}\,|\,f(x)=0\}, \qquad \mcal{R}_{-}(x;P\mcal{K}) =\{(f,g)\in P\mcal{K}\,|\,g(x)=0\}.
\end{gather*}
Then, we can see that $P\mcal{K}_{\{x\},\varnothing}$ is the image of $\mcal{R}_{-}(x;P\mcal{K})$ under the orthogonal projection onto~$\mcal{K}_{\mfrak{X}\backslash\{x\}}$, and similarly, $P\mcal{K}_{\varnothing,\{x\}}$ is the image of $\mcal{R}_{+}(x;P\mcal{K})$ under the orthogonal projection onto~$\mcal{K}_{\mfrak{X}\backslash\{x\}}$.

\subsection[Computation of P(x,0) and P(0,x)]{Computation of $\boldsymbol{P_{\{x\},\varnothing}}$ and $\boldsymbol{P_{\varnothing,\{x\}}}$}

We introduce a lemma found in \cite[Section~7]{BufetovOlshanski2019} that plays prominent roles in computation of~$P_{\{x\},\varnothing}$ and~$P_{\varnothing,\{x\}}$.
\begin{Lemma}[\cite{BufetovOlshanski2019}]\label{lem:projection_reduction_formula}
Let $\mcal{H}$ be a Hilbert space, $\mcal{H}=\mcal{H}_{1}\oplus \mcal{H}_{2}$ be its direct sum decomposition such that $\dim \mcal{H}_{2}<\infty$ and write $\pi_{1}\colon \mcal{H}\to\mcal{H}_{1}$ for the canonical projection. Let $\mcal{L}\subset \mcal{H}$ be a~closed subspace and write the orthogonal projection $P$ onto $\mcal{L}$ in the form of
\begin{equation*}
	P=\left(
	\begin{matrix}
	A & B \\
	C & D
	\end{matrix}
	\right)
\end{equation*}
according to the direct sum decomposition $\mcal{H}=\mcal{H}_{1}\oplus \mcal{H}_{2}$, where $A\in \mfrak{B}(\mcal{H}_{1})$, $B\in \mfrak{B}(\mcal{H}_{2},\mcal{H}_{1})$, $C\in \mfrak{B}(\mcal{H}_{1},\mcal{H}_{2})$, and $D\in \mfrak{B}(\mcal{H}_{2})$.
\begin{enumerate}\itemsep=0pt
\item[$1.$] 	We define a subspace $\mcal{L}_{1}:=\mcal{L}\cap \mcal{H}_{1}$ of $\mcal{H}_{1}$ and write $P_{1}$ for the orthogonal projection onto~$\mcal{L}_{1}$ in~$\mcal{H}_{1}$. If~$D$ is invertible, i.e., $\mcal{L}^{\perp}\cap\mcal{H}_{2}=\{0\}$, we have $P_{1}=A-BD^{-1}C$.
\item[$2.$] 	We define a subspace $\widetilde{\mcal{L}}_{1}:=\pi_{1}(\mcal{L})$ of $\mcal{H}_{1}$ and write $\widetilde{P}_{1}$ for the orthogonal projection onto~$\widetilde{\mcal{L}}_{1}$ in $\mcal{H}_{1}$. If $1-D$ is invertible, i.e., $\mcal{L}\cap\mcal{H}_{2}=\{0\}$, we have $\widetilde{P}_{1}=A+B(1-D)^{-1}C$.
\end{enumerate}
\end{Lemma}
\begin{proof}
Since $P$ is an orthogonal projection, we have $A^{\ast}=A$, $B^{\ast}=C$, $D^{\ast}=D$ and
\begin{gather}\label{eq:relation_component_peojection}
	A^{2}+BC =A, \qquad AB+BD =B, \qquad CA+DC = C, \qquad CB+D^{2} =D.
\end{gather}

1.~Set $\pr{P}_{1}:=A-BD^{-1}C$, which is obviously self-adjoint and is also checked by means of~(\ref{eq:relation_component_peojection}) to be an idempotent.
		Therefore, $\pr{P}_{1}$ is the orthogonal projection onto a closed subspace $\pr{\mcal{L}}_{1}\subset\mcal{H}_{1}$.
		Suppose that $\xi\in\mcal{L}_{1}$. Then $(\xi,0)\in\mcal{H}_{1}\oplus \mcal{H}_{2}$ belongs to $\mcal{L}$ and fixed by $P$:
		\begin{equation*}
		\left(
		\begin{matrix}
		A & B \\
		C & D
		\end{matrix}
		\right)\left(
		\begin{matrix}
		\xi \\
		0
		\end{matrix}
		\right)=\left(
		\begin{matrix}
		\xi \\
		0
		\end{matrix}
		\right),
		\end{equation*}
		which is equivalent to $A\xi=\xi$ and $C\xi=0$.
		Therefore, $\xi$ is fixed by $\pr{P}_{1}$ implying that $\mcal{L}_{1}\subset \pr{\mcal{L}}_{1}$.
		Conversely, suppose that $\xi\in \pr{\mcal{L}}_{1}$. Then, we have $\big(A-BD^{-1}C\big)\xi=\xi$ implying that
		\begin{equation*}
			(\xi,A\xi)-\big(\xi,BD^{-1}C\xi\big)=(\xi,\xi).
		\end{equation*}
		Since $P$ is an orthogonal projection, we have $(\xi,A\xi)\le (\xi,\xi)$. Due to the property $B^{\ast}=C$, the operator $BD^{-1}C$ is a positive operator and $\big(\xi,BD^{-1}C\xi\big)\ge 0$. Therefore, we must have $A\xi=\xi$ and $C\xi=0$ implying $\pr{\mcal{L}}_{1}\subset \mcal{L}_{1}$.

2.~The operator $A+B(1-D)^{-1}C$ is self-adjoint and shown to be an idempotent.
		An element $\xi=(\xi_{1},\xi_{2})\in\mcal{H}_{1}\oplus\mcal{H}_{2}$ belongs to $\mcal{L}$ if and only if
		\begin{equation*}
			A\xi_{1}+B\xi_{2}=\xi_{1},\qquad C\xi_{1}+D\xi_{2}=\xi_{2}.
		\end{equation*}
		Since we have assumed that $1-D$ is invertible, the latter equation gives us $\xi_{2}=(1-D)^{-1}C\xi_{1}$, substitution of which into the former one gives $\big(A+B(1-D)^{-1}C\big)\xi_{1}=\xi_{1}$. Therefore, we obtain the desired result.
\end{proof}

Now, we can obtain explicit expressions of $P_{\{x\},\varnothing}$ and $P_{\varnothing,\{x\}}$.
Let us first write $P\in \mrm{Gr}(\mcal{K},\Gamma)$ as
\begin{equation*}
	P=\left(
	\begin{matrix}
	P_{11} & P_{12} \\
	P_{21} & P_{22}
	\end{matrix}
	\right)
\end{equation*}
according to the direct sum decomposition $\mcal{K}=\ell^{2}(\mfrak{X})\oplus \ell^{2}(\mfrak{X})$, and further write
\begin{equation*}
	P_{ij}=\left(
	\begin{matrix}
		a_{ij} & b_{ij} \\
		c_{ij} & d_{ij}
	\end{matrix}
	\right),\qquad i,j=1,2,
\end{equation*}
according to the direct sum decomposition $\ell^{2}(\mfrak{X})=\ell^{2}(\mfrak{X}\backslash\{x\})\oplus \ell^{2}(\{x\})$.

\begin{Proposition}\label{prop:expression_cond_proj}
Under the above notation, we have, if $P$ is $(\{x\},\varnothing)$-regular,
\begin{equation*}
	P_{\{x\},\varnothing}=\left(
	\begin{matrix}
		a_{11}-b_{12}d_{22}^{-1}c_{21}+b_{11}(1-d_{11})^{-1}c_{11} & a_{12}-b_{12}d_{22}^{-1}c_{22}+b_{11}(1-d_{11})^{-1}c_{12}\\
		a_{21}-b_{22}d_{22}^{-1}c_{21}+b_{21}(1-d_{11})^{-1}c_{11} & a_{22}-b_{22}d_{22}^{-1}c_{22}+b_{21}(1-d_{11})^{-1}c_{12}
	\end{matrix}
	\right)
\end{equation*}
and, if $P$ is $(\varnothing,\{x\})$-regular,
\begin{equation*}
	P_{\varnothing,\{x\}}=\left(
	\begin{matrix}
		a_{11}-b_{11}d_{11}^{-1}c_{11}+b_{12}(1-d_{22})^{-1}c_{21} & a_{12}-b_{11}d_{11}^{-1}c_{12}+b_{12}(1-d_{22})^{-1}c_{22}\\
		a_{21}-b_{21}d_{11}^{-1}c_{11}+b_{22}(1-d_{22})^{-1}c_{21} & a_{22}-b_{21}d_{11}^{-1}c_{12}+b_{22}(1-d_{22})^{-1}c_{22}
	\end{matrix}
	\right)
\end{equation*}
according to the direct sum decomposition $\mcal{K}_{\mfrak{X}\backslash\{x\}}=\mcal{K}^{+}_{\mfrak{X}\backslash\{x\}}\oplus \mcal{K}^{-}_{\mfrak{X}\backslash\{x\}}$.
\end{Proposition}
\begin{proof}
As we saw in Remark \ref{rem:check_anti-sym_matrix_func}, $P_{12}$ and $P_{21}$ are anti-symmetric; in particular, $d_{12}=d_{21}=0$. We also saw that $P_{11}$ and $P_{22}$ are self-adjoint and $JP_{11}J=1-P_{22}$, which implies $d_{11}=1-d_{22}$.

Let us derive the expression of $P_{\{x\},\varnothing}$. By assumption that $P$ is $(\{x\},\varnothing)$-regular, we have $d_{11}\neq 1$ and $d_{22}\neq 0$. Therefore, the desired expression makes sense. It is convenient to write $P$ as
\begin{equation*}
	P=\left(
	\begin{matrix}
	A & b_{1} & b_{2} \\
	c_{1}^{\mrm{T}} & d_{11} & 0 \\
	c_{2}^{\mrm{T}} & 0 & d_{22}
	\end{matrix}
	\right)
\end{equation*}
according to the direct sum decomposition $\mcal{K}=\mcal{K}_{\mfrak{X}\backslash \{x\}}\oplus \mcal{K}_{\{x\}}^{+}\oplus \mcal{K}_{\{x\}}^{-}$, where
\begin{gather*}
	A =\left(
	\begin{matrix}
		a_{11} & a_{12} \\
		a_{21} & a_{22}
	\end{matrix}
	\right), \qquad
	b_{1} =\left(
	\begin{matrix}
		b_{11} \\
		b_{21}
	\end{matrix}
	\right), \qquad
	b_{2} =\left(
	\begin{matrix}
		b_{12} \\
		b_{22}
	\end{matrix}
	\right), \\
	c_{1}^{\mrm{T}} =\left(
	\begin{matrix}
		c_{11} & c_{12}
	\end{matrix}
	\right), \qquad
	c_{2}^{\mrm{T}} =\left(
	\begin{matrix}
		c_{21} & c_{22}
	\end{matrix}
	\right).
\end{gather*}
Notice that $\mcal{R}_{-}(x;\mcal{K})=\mcal{K}_{\mfrak{X}\backslash\{x\}}\oplus \mcal{K}_{\{x\}}^{+}$.
Since $d_{22}\neq 0$, we can apply the formula (1) in Lemma~\ref{lem:projection_reduction_formula} for $\mcal{H}_{1}=\mcal{R}_{-}(x;\mcal{K})$ and $\mcal{H}_{2}=\mcal{K}_{\{x\}}^{-}$ to express the orthogonal projection onto $\mcal{R}_{-}(x;P\mcal{K})$ in $\mcal{R}_{-}(x;\mcal{K})$ as
\begin{equation*}
	\left(
	\begin{matrix}
	A & b_{1} \\
	c_{1}^{\mrm{T}} & d_{11}
	\end{matrix}
	\right)-
	\left(
	\begin{matrix}
		b_{2} \\
		0
	\end{matrix}
	\right)d_{22}^{-1}\left(
	\begin{matrix}
		c_{2}^{\mrm{T}} & 0
	\end{matrix}
	\right)=
	\left(
	\begin{matrix}
		A-b_{2}d_{22}^{-1}c_{2}^{\mrm{T}} & b_{1} \\
		c_{1}^{\mrm{T}} & d_{11}
	\end{matrix}
	\right)
\end{equation*}
according to the direct sum decomposition $\mcal{R}_{-}(x;\mcal{K})=\mcal{K}_{\mfrak{X}\backslash\{x\}}\oplus \mcal{K}_{\{x\}}^{+}$.
Recall that $P\mcal{K}_{\{x\},\varnothing}$ is the image of $\mcal{R}_{-}(x;P\mcal{K})$ under the projection onto $\mcal{K}_{\mfrak{X}\backslash\{x\}}$.
Now, since $d_{11}\neq 1$, we can apply the formula~(2) of Lemma~\ref{lem:projection_reduction_formula} for $\mcal{H}_{1}=\mcal{K}_{\mfrak{X}\backslash\{x\}}$ and $\mcal{H}_{2}=\mcal{K}_{\{x\}}^{+}$ to express $P_{\{x\},\varnothing}$ as
\begin{equation*}
	P_{\{x\},\varnothing}=A-b_{2}d_{22}^{-1}c_{2}^{\mrm{T}}+b_{1}(1-d_{11})^{-1}c_{1}^{\mrm{T}},
\end{equation*}
which is exactly the desired result for $P_{\{x\},\varnothing}$.
The expression of $P_{\varnothing,\{x\}}$ can also be derived in a~similar way to prove the desired result.
\end{proof}

\subsection{Correlation functions of conditional measures}
Let us proceed to computation of the correlation functions under conditional measures.
Let us take $P\in \mrm{Gr}(\mcal{K},\Gamma)$ that is $(\{x\},\varnothing)$-regular. In this case $M^{P}(C(\{x\},\varnothing))=\mbb{K}^{M^{P}}_{12}(x,x)=(e_{x},P_{22}e_{x})$ is non-zero. Therefore, the conditional measure $M^{P}_{\{x\},\varnothing}$ on $\Omega (\mfrak{X}\backslash\{x\})$ is well-defined. Due to the arguments in Section~\ref{subsect:conditioning_CAR_alg}, a correlation function of $M^{P}_{\{x\},\varnothing}$ is computed as
\begin{gather*}
	\rho^{M^{P}_{\{x\},\varnothing}}(x_{1},\dots, x_{n}) =\frac{\varphi_{P}(a_{x_{1}}^{\ast}\cdots a_{x_{n}}^{\ast}a_{x_{n}}\cdots a_{x_{1}}a_{x}^{\ast}a_{x})}{\varphi_{P}(a_{x}^{\ast}a_{x})}
	 =\frac{\pf \left[\mbb{K}_{P}(x_{i},x_{j})\right]_{0\le i,j\le n}}{\mbb{K}^{M^{P}}_{12}(x,x)},
\end{gather*}
where we set $x_{0}:=x$ and $x_{1},\dots, x_{n}\in \mfrak{X}\backslash \{x\}$ are distinct.
It is standard to express the above quantity in terms of the Pfaffian of a matrix of smaller size.
In fact, by means of the formula $\pf (B^{\mrm{T}}AB)=(\det B)(\pf A)$ for an anti-symmetric matrix $A$ and a matrix $B$, we can see that
\begin{equation*}
	\rho^{M^{P}_{\{x\},\varnothing}}(x_{1},\dots, x_{n})=\pf \big[\widetilde{\mbb{K}}(x_{i},x_{j})\big]_{1\le i,j\le n},
\end{equation*}
where
\begin{gather*}
\widetilde{\mbb{K}}_{11}(y,z)=\mbb{K}^{M^{P}}_{11}(y,z)-\frac{\mbb{K}^{M^{P}}_{12}(y,x)\mbb{K}^{M^{P}}_{11}(x,z)}{\mbb{K}^{M^{P}}_{12}(x,x)}+\frac{\mbb{K}^{M^{P}}_{11}(y,x)\mbb{K}^{M^{P}}_{12}(x,z)}{\mbb{K}^{M^{P}}_{12}(x,x)}, \\
\widetilde{\mbb{K}}_{12}(y,z)=\mbb{K}^{M^{P}}_{12}(y,z)-\frac{\mbb{K}^{M^{P}}_{12}(y,x)\mbb{K}^{M^{P}}_{12}(x,z)}{\mbb{K}^{M^{P}}_{12}(x,x)}+\frac{\mbb{K}^{M^{P}}_{11}(y,x)\mbb{K}^{M^{P}}_{22}(x,z)}{\mbb{K}^{M^{P}}_{12}(x,x)}, \\
\widetilde{\mbb{K}}_{22}(y,z)=\mbb{K}^{M^{P}}_{22}(y,z)-\frac{\mbb{K}^{M^{P}}_{22}(y,x)\mbb{K}^{M^{P}}_{12}(x,z)}{\mbb{K}^{M^{P}}_{12}(x,x)}+\frac{\mbb{K}^{M^{P}}_{12}(y,x)\mbb{K}^{M^{P}}_{22}(x,z)}{\mbb{K}^{M^{P}}_{12}(x,x)},
\end{gather*}
for $y,z\in\mfrak{X}\backslash\{x\}$. Note that the other component $\widetilde{\mbb{K}}_{21}(y,z)$ is determined by the anti-symmetry $\widetilde{\mbb{K}}(y,z)^{\mrm{T}}=-\widetilde{\mbb{K}}(z,y)$.

We next assume that $P\in \mrm{Gr}(\mcal{K},\Gamma)$ is $(\varnothing,\{x\})$-regular. Then, $M^{P}(C(\varnothing,\{x\}))=1-\mbb{K}^{M^{P}}_{12}(x,x)=(e_{x},P_{11}e_{y})$ is non-zero. Therefore, the conditional measure $M^{P}_{\varnothing,\{x\}}$ on $\Omega (\mfrak{X}\backslash\{x\})$ is well-defined. A correlation function is computed as
\begin{gather*}
	\rho^{M^{P}_{\varnothing,\{x\}}}(x_{1},\dots, x_{n}) =\frac{\varphi_{P}(a_{x_{1}}^{\ast}\cdots a_{x_{n}}^{\ast}a_{x_{n}}\cdots a_{x_{1}}a_{x}a^{\ast}_{x})}{\varphi_{P}(a_{x}a^{\ast}_{x})}
	 =\frac{\pf [\pr{\mbb{K}}_{P}(x_{i},x_{j}) ]_{0\le i,j\le n}}{1-\mbb{K}^{M^{P}}_{12}(x,x)},
\end{gather*}
where we set $x_{0}:=x$, $x_{1},\dots, x_{n}\in \mfrak{X}\backslash\{x\}$ are distinct and
\begin{gather*}
	\pr{\mbb{K}}_{P}(x_{i},x_{j}):=
	\begin{cases}
		\mbb{K}_{P}(x_{i},x_{j}), & 1\le i,j\le n, \\
		\left(
		\begin{matrix}
			(e_{x},P_{11}e_{x_{j}}) & (e_{x},P_{12}e_{x_{j}}) \\
			(e_{x}, P_{21}e_{x_{j}}) & (e_{x},P_{22}e_{x_{j}})
		\end{matrix}
		\right), & i=0,\quad 1\le j\le n, \vspace{1mm}\\
		\left(
		\begin{matrix}
			(e_{x_{i}},P_{22}e_{x}) & (e_{x_{i}},P_{21}e_{x}) \\
			(e_{x_{i}}, P_{12}e_{x}) & (e_{x_{i}},P_{11}e_{x})
		\end{matrix}
		\right), & 1\le i \le n,\quad j=0, \vspace{1mm}\\
		\left(
		\begin{matrix}
			0 & (e_{x},P_{11}e_{x}) \\
			(e_{x}, (P_{22}-1)e_{x}) & 0
		\end{matrix}
		\right), & i=j=0.
	\end{cases}
\end{gather*}
Again, we write the correlation function in terms of the Pfaffian of a $(2n\times 2n)$-anti-symmetric matrix. Consequently, we obtain
\begin{equation*}
	\rho^{M^{P}_{\varnothing,\{x\}}}(x_{1},\dots, x_{n})=\pf\big[\what{\mbb{K}}(x_{i},x_{j})\big]_{1\le i,j\le n},
\end{equation*}
where
\begin{gather*}
\what{\mbb{K}}_{11}(y,z) =\mbb{K}^{M^{P}}_{11}(y,z)-\frac{\mbb{K}^{M^{P}}_{11}(y,x)\mbb{K}^{M^{P}}_{12}(x,z)}{1-\mbb{K}^{M^{P}}_{12}(x,x)}+\frac{\mbb{K}^{M^{P}}_{12}(y,x)\mbb{K}^{M^{P}}_{11}(x,z)}{1-\mbb{K}^{M^{P}}_{12}(x,x)}, \\
\what{\mbb{K}}_{12}(y,z) =\mbb{K}^{M^{P}}_{12}(y,z)-\frac{\mbb{K}^{M^{P}}_{11}(y,x)\mbb{K}^{M^{P}}_{22}(x,z)}{1-\mbb{K}^{M^{P}}_{12}(x,x)}+\frac{\mbb{K}^{M^{P}}_{12}(y,x)\mbb{K}^{M^{P}}_{12}(x,z)}{1-\mbb{K}^{M^{P}}_{12}(x,x)}, \\
\what{\mbb{K}}_{22}(y,z) =\mbb{K}^{M^{P}}_{22}(y,z)-\frac{\mbb{K}^{M^{P}}_{12}(y,x)\mbb{K}^{M^{P}}_{22}(x,z)}{1-\mbb{K}^{M^{P}}_{12}(x,x)}+\frac{\mbb{K}^{M^{P}}_{22}(y,x)\mbb{K}^{M^{P}}_{12}(x,z)}{1-\mbb{K}^{M^{P}}_{12}(x,x)}
\end{gather*}
for $y,z\in\mfrak{X}\backslash\{x\}$, and the other component is determined by the anti-symmetry $\what{\mbb{K}}(y,z)^{\mrm{T}}=-\what{\mbb{K}}(z,y)$.

\subsection{Proof of Theorem \ref{thm:conditional_projection}}
It remains to show that $\widetilde{\mbb{K}}(y,z)=\mbb{K}_{P_{\{x\},\varnothing}}(y,z)$ and $\what{\mbb{K}}(y,z)=\mbb{K}_{P_{\varnothing,\{x\}}}(y,z)$, $y,z\in \mfrak{X}\backslash\{x\}$.
It is straightforward that
\begin{gather*}
\widetilde{\mbb{K}}_{11}(y,z)=(e_{y},P_{21}e_{z})-\frac{(e_{y},P_{22}e_{x})(e_{x},P_{21}e_{z})}{(e_{x},P_{22}e_{x})}+\frac{(e_{y},P_{21}e_{x})(e_{x},P_{11}e_{z})}{(e_{x},(1-P_{11})e_{x})}, \\
\widetilde{\mbb{K}}_{12}(y,z)=(e_{y},P_{22}e_{z})-\frac{(e_{y},P_{22}e_{x})(e_{x},P_{22}e_{z})}{(e_{x},P_{22}e_{x})}+\frac{(e_{y},P_{21}e_{x})(e_{x},P_{12}e_{z})}{(e_{x},(1-P_{11})e_{x})}, \\
\widetilde{\mbb{K}}_{22}(y,z)=(e_{y},P_{12}e_{z})-\frac{(e_{y},P_{12}e_{x})(e_{x},P_{22}e_{z})}{(e_{x},P_{22}e_{x})}+\frac{(e_{y},P_{11}e_{x})(e_{x},P_{12}e_{z})}{(e_{x},(1-P_{11})e_{x})}.
\end{gather*}
Comparing these descriptions with the result of Proposition~\ref{prop:expression_cond_proj}, we conclude that $\widetilde{\mbb{K}}(y,z)=\mbb{K}_{P_{\{x\},\varnothing}}(y,z)$, $y,z\in\mfrak{X}\backslash\{x\}$.
It also follows from
\begin{gather*}
\what{\mbb{K}}_{11}(y,z)=(e_{y},P_{21}e_{z})-\frac{(e_{y}P_{21}e_{x})(e_{x},P_{11}e_{z})}{(e_{x},P_{11}e_{x})}+\frac{(e_{y},P_{22}e_{x})(e_{x},P_{21}e_{z})}{(e_{x},(1-P_{22})e_{x})}, \\
\what{\mbb{K}}_{12}(y,z)=(e_{y},P_{22}e_{z})-\frac{(e_{y}P_{21}e_{x})(e_{x},P_{12}e_{z})}{(e_{x},P_{11}e_{x})}+\frac{(e_{y},P_{22}e_{x})(e_{x},P_{22}e_{z})}{(e_{x},(1-P_{22})e_{x})}, \\
\what{\mbb{K}}_{22}(y,z)=(e_{y},P_{12}e_{z})-\frac{(e_{y}P_{11}e_{x})(e_{x},P_{12}e_{z})}{(e_{x},P_{11}e_{x})}+\frac{(e_{y},P_{12}e_{x})(e_{x},P_{22}e_{z})}{(e_{x},(1-P_{22})e_{x})}
\end{gather*}
and Proposition \ref{prop:expression_cond_proj} that the other desired coincidence $\what{\mbb{K}}(y,z)=\mbb{K}_{P_{\varnothing,\{x\}}}(y,z)$, $y,z\in \mfrak{X}\backslash\{x\}$ holds.

\begin{thebibliography}{99}
\bibitem[2]{Araki1970}
Araki H., On quasifree states of {${\rm CAR}$} and {B}ogoliubov automorphisms,
 \href{https://doi.org/10.2977/prims/1195193913}{\textit{Publ. Res. Inst. Math. Sci.}} \textbf{6} (1970), 385--442.

\bibitem[3]{Binnenhei1995}
Binnenhei C., Implementation of endomorphisms of the {CAR} algebra,
 \href{https://doi.org/10.1142/S0129055X95000323}{\textit{Rev. Math. Phys.}} \textbf{7} (1995), 833--869,
 \href{https://arxiv.org/abs/physics/9701009}{arXiv:physics/9701009}.

\bibitem[5]{Borodin2011}
Borodin A., Determinantal point processes, in The {O}xford Handbook of Random
 Matrix Theory, Oxford University Press, Oxford, 2011, 231--249,
 \href{https://arxiv.org/abs/0911.1153}{arXiv:0911.1153}.

\bibitem[6]{Borodin1998a}
Borodin A., Olshanski G., Point processes and the infinite symmetric group.
 Part~II: Higher correlation functions, \href{https://arxiv.org/abs/math.RT/9804087}{arXiv:math.RT/9804087}.

\bibitem[7]{BorodinOlshanski1998a}
Borodin A., Olshanski G., Point processes and the infinite symmetric group.
 Part~III: Fermion point processes, \href{https://arxiv.org/abs/math.RT/9804088}{arXiv:math.RT/9804088}.

\bibitem[8]{Borodin1998b}
Borodin A., Olshanski G., Point processes and the infinite symmetric group.
 Part~IV: Matrix Whittaker kernel, \href{https://arxiv.org/abs/math.RT/9810013}{arXiv:math.RT/9810013}.

\bibitem[9]{BorodinOlshanski1998b}
Borodin A., Olshanski G., Point processes and the infinite symmetric group,
 \href{https://doi.org/10.4310/MRL.1998.v5.n6.a9}{\textit{Math. Res. Lett.}} \textbf{5} (1998), 799--816,
 \href{https://arxiv.org/abs/math.RT/9810015}{arXiv:math.RT/9810015}.

\bibitem[10]{BorodinRains2005}
Borodin A., Rains E.M., Eynard--{M}ehta theorem, {S}chur process, and their
 {P}faffian analogs, \href{https://doi.org/10.1007/s10955-005-7583-z}{\textit{J.~Stat. Phys.}} \textbf{121} (2005), 291--317,
 \href{https://arxiv.org/abs/math-ph/0409059}{arXiv:math-ph/0409059}.

\bibitem[11]{BratteliRobinson1987}
Bratteli O., Robinson D.W., Operator algebras and quantum statistical
 mechanics. 1.~$C^\ast$- and $W^\ast$-algebras, symmetry groups, decomposition
 of states, 2nd~ed., \textit{Texts and Monographs in Physics}, \href{https://doi.org/10.1007/978-3-662-02520-8}{Springer-Verlag}, New
 York, 1987.

\bibitem[12]{BratteliRobinson1997}
Bratteli O., Robinson D.W., Operator algebras and quantum statistical
 mechanics. 2.~Equilibrium states. Models in quantum statistical mechanics,
 2nd~ed., \textit{Texts and Monographs in Physics}, \href{https://doi.org/10.1007/978-3-662-03444-6}{Springer-Verlag}, Berlin, 1997.

\bibitem[13]{BufetovCundenQiu2019}
Bufetov A.I., Cunden F.D., Qiu Y., {$J$}-{H}ermitian determinantal point
 processes: balanced rigidity and balanced {P}alm equivalence
 \href{https://arxiv.org/abs/1912.10743}{arXiv:1912.10743}.

\bibitem[14]{BufetovOlshanski2019}
Bufetov A.I., Olshanski G., A hierarchy of {P}alm measures for determinantal
 point processes with gamma kernels, \href{https://arxiv.org/abs/1904.13371}{arXiv:1904.13371}.

\bibitem[16]{Ferrari2004}
Ferrari P.L., Polynuclear growth on a flat substrate and edge scaling of {GOE}
 eigenvalues, \href{https://doi.org/10.1007/s00220-004-1204-6}{\textit{Comm. Math. Phys.}} \textbf{252} (2004), 77--109,
 \href{https://arxiv.org/abs/math-ph/0402053}{arXiv:math-ph/0402053}.

\bibitem[17]{BenHoughKrishnapurPeresVirag2006}
Hough J.B., Krishnapur M., Peres Y., Vir\'ag B., Determinantal processes and
 independence, \href{https://doi.org/10.1214/154957806000000078}{\textit{Probab. Surv.}} \textbf{3} (2006), 206--229,
 \href{https://arxiv.org/abs/math.PR/0503110}{arXiv:math.PR/0503110}.

\bibitem[21]{KatoriTanemura2007}
Katori M., Tanemura H., Infinite systems of noncolliding generalized meanders
 and {R}iemann--{L}iouville differintegrals, \href{https://doi.org/10.1007/s00440-006-0015-4}{\textit{Probab. Theory Related
 Fields}} \textbf{138} (2007), 113--156, \href{https://arxiv.org/abs/math.PR/0506187}{arXiv:math.PR/0506187}.

\bibitem[23]{Koshmanenko1999}
Koshmanenko V., Singular quadratic forms in perturbation theory,
 \textit{Mathematics and its Applications}, Vol.~474, \href{https://doi.org/10.1007/978-94-011-4619-7}{Kluwer Academic
 Publishers}, Dordrecht, 1999.

\bibitem[24]{Lenard1975a}
Lenard A., States of classical statistical mechanical systems of infinitely
 many particles.~{I}, \href{https://doi.org/10.1007/BF00251601}{\textit{Arch. Rational Mech. Anal.}} \textbf{59} (1975),
 219--239.

\bibitem[25]{Lenard1975b}
Lenard A., States of classical statistical mechanical systems of infinitely
 many particles. {II}.~Characterization of correlation measures, \href{https://doi.org/10.1007/BF00251602}{\textit{Arch.
 Rational Mech. Anal.}} \textbf{59} (1975), 240--256.

\bibitem[26]{Lyons2003}
Lyons R., Determinantal probability measures, \href{https://doi.org/10.1007/s10240-003-0016-0}{\textit{Publ. Math. Inst. Hautes
 \'Etudes Sci.}} (2003), 167--212, \href{https://arxiv.org/abs/math.PR/0204325}{arXiv:math.PR/0204325}.

\bibitem[27]{Lytvynov2002}
Lytvynov E., Fermion and boson random point processes as particle distributions
 of infinite free {F}ermi and {B}ose gases of finite density, \href{https://doi.org/10.1142/S0129055X02001533}{\textit{Rev.
 Math. Phys.}} \textbf{14} (2002), 1073--1098, \href{https://arxiv.org/abs/math-ph/0112006}{arXiv:math-ph/0112006}.

\bibitem[28]{LytvynovMei2007}
Lytvynov E., Mei L., On the correlation measure of a family of commuting
 {H}ermitian operators with applications to particle densities of the
 quasi-free representations of the {CAR} and {CCR}, \href{https://doi.org/10.1016/j.jfa.2006.12.017}{\textit{J.~Funct. Anal.}}
 \textbf{245} (2007), 62--88, \href{https://arxiv.org/abs/math.PR/0608334}{arXiv:math.PR/0608334}.

\bibitem[31]{Matsumoto2005}
Matsumoto S., Correlation functions of the shifted {S}chur measure,
 \href{https://doi.org/10.2969/jmsj/1158241925}{\textit{J.~Math. Soc. Japan}} \textbf{57} (2005), 619--637,
 \href{https://arxiv.org/abs/math.CO/0312373}{arXiv:math.CO/0312373}.

\bibitem[32]{MatsumotoShirai2013}
Matsumoto S., Shirai T., Correlation functions for zeros of a {G}aussian power
 series and {P}faffians, \href{https://doi.org/10.1214/EJP.v18-2545}{\textit{Electron.~J. Probab.}} \textbf{18} (2013),
 no.~49, 18~pages, \href{https://arxiv.org/abs/1212.6108}{arXiv:1212.6108}.

\bibitem[33]{Nagao2007}
Nagao T., Pfaffian expressions for random matrix correlation functions,
 \href{https://doi.org/10.1007/s10955-007-9415-9}{\textit{J.~Stat. Phys.}} \textbf{129} (2007), 1137--1158, \href{https://arxiv.org/abs/0708.2036}{arXiv:0708.2036}.

\bibitem[35]{Okounkov2001}
Okounkov A., Infinite wedge and random partitions, \href{https://doi.org/10.1007/PL00001398}{\textit{Selecta Math.
 (N.S.)}} \textbf{7} (2001), 57--81, \href{https://arxiv.org/abs/math.RT/9907127}{arXiv:math.RT/9907127}.

\bibitem[36]{Olshanski1998}
Olshanski G., Point processes and the infinite symmetric group. Part~V:
 Analysis of the matrix Whittaker kernel, \href{https://arxiv.org/abs/math.RT/9810014}{arXiv:math.RT/9810014}.

\bibitem[37]{Olshanski2003}
Olshanski G., Point processes related to the infinite symmetric group, in The
 Orbit Method in Geometry and Physics ({M}arseille, 2000), \textit{Progr.
 Math.}, Vol.~213, \href{https://doi.org/10.1007/978-1-4612-0029-1_15}{Birkh\"auser Boston}, Boston, MA, 2003, 349--393.

\bibitem[39]{Olshanski2020}
Olshanski G., Determinantal point processes and fermion quasifree states,
 \href{https://doi.org/10.1007/s00220-020-03716-1}{\textit{Comm. Math. Phys.}} \textbf{378} (2020), 507--555,
 \href{https://arxiv.org/abs/2002.10723}{arXiv:2002.10723}.

\bibitem[43]{Rains2000}
Rains E.M., Correlation functions for symmetrized increasing subsequences,
 \href{https://arxiv.org/abs/math.CO/0006097}{arXiv:math.CO/0006097}.

\bibitem[44]{RedeiSummers2006}
R\'edei M., Summers S.J., Quantum probability theory, \href{https://doi.org/10.1016/j.shpsb.2006.05.006}{\textit{Stud. Hist.
 Philos. Sci.~B Stud. Hist. Philos. Modern Phys.}} \textbf{38} (2007),
 390--417, \href{https://arxiv.org/abs/quant-ph/0601158}{arXiv:quant-ph/0601158}.

\bibitem[47]{ShiraiTakahashi2003a}
Shirai T., Takahashi Y., Random point fields associated with certain {F}redholm
 determinants. {I}.~{F}ermion, {P}oisson and boson point processes,
 \href{https://doi.org/10.1016/S0022-1236(03)00171-X}{\textit{J.~Funct. Anal.}} \textbf{205} (2003), 414--463.

\bibitem[48]{ShiraiTakahashi2003b}
Shirai T., Takahashi Y., Random point fields associated with certain {F}redholm
 determinants. {II}.~{F}ermion shifts and their ergodic and {G}ibbs
 properties, \href{https://doi.org/10.1214/aop/1055425789}{\textit{Ann. Probab.}} \textbf{31} (2003), 1533--1564.

\bibitem[49]{Soshnikov2000}
Soshnikov A., Determinantal random point fields, \href{https://doi.org/10.1070/rm2000v055n05ABEH000321}{\textit{Russian Math. Surveys}}
 \textbf{55} (2000), 923--975.

\bibitem[50]{Spohn1987}
Spohn H., Interacting {B}rownian particles: a study of {D}yson's model, in
 Hydrodynamic Behavior and Interacting Particle Systems ({M}inneapolis,
 {M}inn., 1986), \textit{IMA Vol. Math. Appl.}, Vol.~9, \href{https://doi.org/10.1007/978-1-4684-6347-7_13}{Springer}, New York,
 1987, 151--179.

\bibitem[51]{Takesaki1979}
Takesaki M., Theory of operator algebras.~{I}, \href{https://doi.org/10.1007/978-1-4612-6188-9}{Springer-Verlag}, New York~--
 Heidelberg, 1979.

\bibitem[54]{Vuletic2007}
Vuleti\'c M., The shifted {S}chur process and asymptotics of large random
 strict plane partitions, \href{https://doi.org/10.1093/imrn/rnm043}{\textit{Int. Math. Res. Not.}} \textbf{2007} (2007),
 rnm043, 53~pages, \href{https://arxiv.org/abs/math-ph/0702068}{arXiv:math-ph/0702068}.

\bibitem[55]{WangLi2019}
Wang Z.-L., Li S.-H., B{KP} hierarchy and {P}faffian point process,
 \href{https://doi.org/10.1016/j.nuclphysb.2018.12.028}{\textit{Nuclear Phys.~B}} \textbf{939} (2019), 447--464, \href{https://arxiv.org/abs/1807.02259}{arXiv:1807.02259}.

\end{thebibliography}
\end{document}
